\chapter{Complexity — Code Simulator}
%%% generated by the blueprint pipeline from TCSlib.Complexity.TuringMachine.CodeSimulator
\section{Overview}

Declarations of the Lean module \texttt{TCSlib.Complexity.TuringMachine.CodeSimulator}.

\section{Declarations}

\begin{definition}[Serialization of a multitape machine]
\lean{Turing.serializeK}
\label{Turing.serializeK}
\leanok
\uses{Turing.MultiTapeTM, Turing.pairEncode, Turing.tableBitsK, Turing.unaryFin}
Let $M$ be a deterministic machine with $k$ work tapes, Boolean symbols, and state set
of size $n+1$. Its serialization is the self-delimiting pairing of the bit
representation of $n$ with the unary encoding of its initial state followed by the bit
encoding of its transition table.
\end{definition}

\begin{definition}[Repeated table skipper]
\lean{Turing.codeSkipTableK}
\label{Turing.codeSkipTableK}
\leanok
\uses{Turing.codeSkipActionK, Turing.codeSkipRepeat}
For natural numbers $k$ and $n$, define a partial function on finite bit strings by
taking the action skipper for $k$ work tapes and $n+1$ states, iterating that skipper
$k+1$ times with each application itself repeated three times, and then repeating this
entire procedure $n+1$ times. The function returns the remaining suffix when all
applications succeed, and is undefined otherwise.
\end{definition}

\begin{lemma}[Consing a value onto a one-entry function]
\lean{Turing.fin_cons_one}
\label{Turing.fin_cons_one}
\leanok
\difficulty{4}
For every optional Boolean value $w$, the function from the one-element finite type to
optional Booleans obtained by placing $w$ first in an empty function is the constant
function with value $w$.
\end{lemma}

\begin{lemma}[Consing values onto a two-entry function]
\lean{Turing.fin_cons_two}
\label{Turing.fin_cons_two}
\leanok
\uses{Turing.workPair}
\difficulty{4}
For any optional Boolean values $w_0$ and $w_1$, the function from the two-element
finite type obtained by placing $w_0$ first and then using the constant one-entry
function with value $w_1$ agrees with the function whose value at index $0$ is $w_0$ and
whose value at index $1$ is $w_1$.
\end{lemma}

\begin{lemma}[Read patterns for one work tape]
\lean{Turing.workReadsK_one}
\label{Turing.workReadsK_one}
\leanok
\uses{Turing.fin_cons_one, Turing.workReadsK}
\difficulty{3}
The list of all read patterns for one Boolean work tape is exactly the list of constant
functions whose value is blank, false, or true.
\end{lemma}

\begin{lemma}[Read patterns for two work tapes]
\lean{Turing.workReadsK_two}
\label{Turing.workReadsK_two}
\leanok
\uses{Turing.fin_cons_one, Turing.fin_cons_two, Turing.workPair, Turing.workReadsK}
\difficulty{3}
The list of all read patterns for two Boolean work tapes is exactly the nested list
obtained by choosing independently, for each tape, a value from blank, false, and true,
and assigning the first choice to tape $0$ and the second to tape $1$.
\end{lemma}

\begin{lemma}[Two-tape table serialization]
\lean{Turing.tableBitsK_two}
\label{Turing.tableBitsK_two}
\leanok
\uses{Turing.Action, Turing.codeTable2Bits, Turing.tableBitsK, Turing.workPair, Turing.workReadsK_two}
\difficulty{3}
Let $n$ be a natural number, and let $\mathrm{tr}$ assign an action on two Boolean work
tapes to each state in $\{0,\ldots,n\}$, each optional Boolean input symbol, and each
pair of optional Boolean symbols read on the work tapes. Serialize the complete
transition table by listing, for every state and every such input and pair of work-tape
reads, the action bits prescribed by $\mathrm{tr}$. This list is exactly the two-tape
table encoding obtained by supplying the two work-tape symbols separately to
$\mathrm{tr}$.
\end{lemma}

\begin{lemma}[Equivalence of two table readers]
\lean{Turing.codeSkipTable2_eq_codeSkipTableK}
\label{Turing.codeSkipTable2_eq_codeSkipTableK}
\leanok
\uses{Turing.codeSkipTable2, Turing.codeSkipTableK}
\difficulty{1}
For every natural number $n$ and Boolean string $xs$, the reader that parses a table
using two work-field repetitions per action equals the reader that parses it using the
iterated action reader with the corresponding two-work-field setting.
\end{lemma}

\begin{definition}[Machine format]
\lean{Turing.MachineFormat}
\label{Turing.MachineFormat}
\leanok
\uses{Turing.CodeFormat, Turing.MultiTapeTM, Turing.codeSkipTableK, Turing.serializeK}
A machine format consists of a code type, a family of transition-table types indexed by
the number of states minus one, a state-count function, and a way to assemble a code
from a state count, an initial state, and a table. It also specifies serializations of
codes and tables, table-reading and table-skipping procedures, a guard constant, and a
fallback code. Assembly records the stated state count, serializes as the pairing of the
binary expansion of that count with the unary initial state followed by the table bits,
and represents every code; table reading correctly parses and reconstructs the table and
suffix, its erased version agrees with table skipping, and every table has at least the
guard constant times the number of states in bits.
\end{definition}

\begin{definition}[Finite machine associated with a code]
\lean{Turing.MachineFormat.toFinTM}
\label{Turing.MachineFormat.toFinTM}
\leanok
\uses{Turing.FinTM, Turing.MachineFormat}
Let a machine format have a code set, assign each code $M$ a natural number $n_M$ and a
binary multi-tape Turing machine with $k$ work tapes and state set $\{0,\ldots,n_M\}$.
The code $M$ determines the bundled finite binary Turing machine with those $k$ work
tapes, that state set, and the assigned machine.
\end{definition}

\begin{lemma}[Agreement of one-tape table skippers]
\lean{Turing.oneTape_skipTable_eq}
\label{Turing.oneTape_skipTable_eq}
\leanok
\uses{Turing.codeSkipAction, Turing.codeSkipActionK, Turing.codeSkipAction_eq_codeSkipActionK, Turing.codeSkipRepeat, Turing.codeSkipTableK, Turing.oneTapeFormat}
\difficulty{3}
For every natural number $n$ and every binary word $x$, the one-tape-format table
skipper with bound $n$ returns the same result as the general table skipper with one
work tape and bound $n$.
\end{lemma}

\begin{lemma}[Agreement of one-tape serializations]
\lean{Turing.oneTape_serialize_eq}
\label{Turing.oneTape_serialize_eq}
\leanok
\uses{Turing.CodeTM, Turing.CodeTM.serialize, Turing.actionBits_eq_actionBitsK, Turing.serializeK, Turing.tableBitsK, Turing.workReadsK_one}
\difficulty{3}
For every machine $M$ in code normal form, its canonical serialization agrees with the
$k$-work-tape serialization of its underlying machine specialized to one work tape.
\end{lemma}

\begin{definition}[One-tape machine format]
\lean{Turing.oneTapeMachineFormat}
\label{Turing.oneTapeMachineFormat}
\leanok
\uses{Turing.MachineFormat, Turing.oneTapeFormat, Turing.oneTape_serialize_eq, Turing.oneTape_skipTable_eq}
Define a machine format by specializing the one-tape code format to one work tape: use
each code's underlying machine as its denoted machine, and use the one-tape
serialization and table-skipping procedures for the required format operations.
\end{definition}

\begin{lemma}[Repeated skipper consumption bound]
\lean{Turing.MachineFormat.skipRepeat_consumes}
\label{Turing.MachineFormat.skipRepeat_consumes}
\leanok
\uses{Turing.codeSkipRepeat}
\difficulty{4}
Let $r$ be a partial reader on finite bit strings and let $b$ be a natural number.
Suppose every successful read from $x$ to suffix $y$ satisfies $b+|y|\le |x|$. If
applying the reader $m$ times to $x$ succeeds with final suffix $y$, then $bm+|y|\le
|x|$.
\end{lemma}

\begin{lemma}[Iterated triple-skip consumption bound]
\lean{Turing.MachineFormat.skipSymbols_consumes}
\label{Turing.MachineFormat.skipSymbols_consumes}
\leanok
\uses{Turing.MachineFormat.skipRepeat_consumes, Turing.codeSkipRepeat}
\difficulty{4}
Let $r$ be a partial reader on finite bit strings and let $b$ be a natural number.
Suppose every successful read from $x$ to suffix $y$ satisfies $b+|y|\le |x|$. If the
reader obtained by iterating the operation $f\mapsto (x\mapsto\text{apply }f\text{ three
times})$ exactly $j$ times, starting from $r$, succeeds on $x$ with suffix $y$, then
$b3^j+|y|\le |x|$.
\end{lemma}

\begin{lemma}[Skip table length bound]
\lean{Turing.MachineFormat.skipTable_consumes}
\label{Turing.MachineFormat.skipTable_consumes}
\leanok
\uses{Turing.MachineFormat.skipRepeat_consumes, Turing.MachineFormat.skipSymbols_consumes, Turing.codeActionK_length, Turing.codeEraseActionK, Turing.codeReadActionK_sound, Turing.codeSkipActionK, Turing.codeSkipTableK}
\difficulty{4}
For natural numbers $k,n$ and bit strings $xs,ys$, if the table skipper with parameters
$k,n$ accepts $xs$ and returns $ys$, then $(4k+5)3^{k+1}(n+1)+|ys|\le |xs|$.
\end{lemma}

\begin{lemma}[Scan grammar length bound]
\lean{Turing.MachineFormat.scan_grammar_guard}
\label{Turing.MachineFormat.scan_grammar_guard}
\leanok
\uses{Turing.CodeFormat.scan, Turing.MachineFormat.skipTable_consumes, Turing.codeBitsNat, Turing.codeEraseFin, Turing.codeReadFin_sound, Turing.pairDecode}
\difficulty{5}
Let $F$ be a machine format, and let $\alpha$, $b$, $r$, and $u$ be finite bit strings.
If parsing $\alpha$ yields the pair $(b,r)$, and scanning $\alpha$ with format $F$
succeeds with remaining suffix $u$, then $(4k+5)3^{k+1}(\operatorname{val}(b)+1)\le
|r|$, where $k$ is the parameter of $F$ and $\operatorname{val}(b)$ is the natural
number represented by $b$ in least-significant-bit-first binary.
\end{lemma}

\begin{definition}[Length-guarded machine scan]
\lean{Turing.MachineFormat.guardedScan}
\label{Turing.MachineFormat.guardedScan}
\leanok
\uses{Turing.CodeFormat.scan, Turing.codeBitsNat, Turing.pairDecode}
For a machine format $F$, define a partial function on finite bit strings as follows.
Decode the input using the pair decoder, obtaining a bit string $b$ and a remainder $r$;
if $r$ has length less than $(4k+5)3^{k+1}(
u(b)+1)$, where $k$ is the parameter of $F$ and $
u(b)$ is the natural-number code length of $b$, return no value, and otherwise return
the result of applying $F$’s scan operation to the original input.
\end{definition}

\begin{lemma}[Guarded scan agrees with scan]
\lean{Turing.MachineFormat.guardedScan_eq}
\label{Turing.MachineFormat.guardedScan_eq}
\leanok
\uses{Turing.CodeFormat.scan, Turing.MachineFormat.guardedScan, Turing.MachineFormat.scan_grammar_guard, Turing.pairDecode}
\difficulty{4}
For every machine format $F$ and every finite Boolean string $\alpha$, the guarded scan
of $\alpha$ equals the scan of $\alpha$ in format $F$.
\end{lemma}

\begin{lemma}[Fallback below the grammar length bound]
\lean{Turing.MachineFormat.decode_of_grammar_guard}
\label{Turing.MachineFormat.decode_of_grammar_guard}
\leanok
\uses{Turing.CodeFormat.decode, Turing.CodeFormat.decode_of_parseFull_none, Turing.CodeFormat.parseFull, Turing.CodeFormat.scan, Turing.CodeFormat.scan_full, Turing.MachineFormat.guardedScan, Turing.MachineFormat.guardedScan_eq, Turing.codeBitsNat, Turing.pairDecode}
\difficulty{3}
Let $F$ be a machine format, and suppose a Boolean string $\alpha$ parses as a paired
header $bits$ and suffix $rest$. If $|rest| < (4k+5)3^{k+1}(v+1)$, where $k$ is the
number of work tapes in $F$ and $v$ is the natural number represented by $bits$ in
least-significant-bit-first binary, then decoding $\alpha$ returns the fallback code of
$F$.
\end{lemma}

\begin{lemma}[A lower bound on the state count]
\lean{Turing.MachineFormat.grammar_guard_states}
\label{Turing.MachineFormat.grammar_guard_states}
\leanok
\difficulty{3}
Let $F$ be a machine format, and let $n,s\in\mathbb N$. If $(4k+5)3^{k+1}(n+1)\le s$,
where $k$ is the parameter of $F$, then $n+1\le s$.
\end{lemma}

\begin{definition}[Canonical decoded header]
\lean{Turing.MachineFormat.canonHeader}
\label{Turing.MachineFormat.canonHeader}
\leanok
\uses{Turing.pairDecode}
For a machine format $F$ and a bit string $\alpha$, define its canonical header to be
the header returned by decoding $\alpha$ as a pair, or the empty bit string if decoding
fails.
\end{definition}

\begin{lemma}[Length bound for the canonical header]
\lean{Turing.MachineFormat.canonHeader_length}
\label{Turing.MachineFormat.canonHeader_length}
\leanok
\uses{Turing.MachineFormat.canonHeader, Turing.eq_pairEncode_of_pairDecode, Turing.length_pairEncode, Turing.pairDecode}
\difficulty{4}
For every Boolean word $\alpha$, the length of its canonical header is at most
$|\alpha|$.
\end{lemma}

\begin{lemma}[Canonical Header Call Machine]
\lean{Turing.MachineFormat.exists_canonHeaderCall}
\label{Turing.MachineFormat.exists_canonHeaderCall}
\leanok
\uses{Turing.Cfg.ofWords, Turing.FinTM, Turing.FinTM.computesFunInTime_pairFst, Turing.FinTM.exists_inputCallTM, Turing.MachineFormat.canonHeader, Turing.MachineFormat.canonHeader_length, Turing.MultiTapeTM.runFrom, Turing.stateWord}
\difficulty{4}
For every machine format $F$, there exist a finite Turing machine $C$ over the Boolean
alphabet with at least one work tape, entry and exit states, and a constant
$a\in\mathbb N$. For every Boolean input $\alpha$, there is a time $t$ with $0<t\le
a(|\alpha|+1)$ such that the configuration at time $t$ has the exit state, the
canonical header of $\alpha$ on tape zero, and empty words on the other work tapes,
with input head at position $1$ and empty output; at each intermediate time $0<u<t$ the
state is not the exit state.
\end{lemma}

\begin{definition}[Guard phase]
\lean{Turing.MachineFormat.guardPhase}
\label{Turing.MachineFormat.guardPhase}
\leanok
For a machine format $F$ and natural numbers $b$ and $u$, define the guard phase to be
the residue of $u$ modulo $b+1$, regarded as an element of $\{0,\ldots,b\}$.
\end{definition}

\begin{definition}[Guard block machine]
\lean{Turing.MachineFormat.guardBlockTM}
\label{Turing.MachineFormat.guardBlockTM}
\leanok
\uses{Turing.MachineFormat.guardPhase, Turing.MultiTapeTM, Turing.inputMoveAction}
For a machine format $F$ and a natural number $b$, define a Turing machine with no work
tapes, binary input, and state set consisting of a halt state together with phases $0$
through $b$. It starts at phase $0$; in the halt state it halts, and at any phase it
halts if the scanned input is blank, while on a present input symbol it moves the input
head one step and advances the phase by one modulo $b+1$.
\end{definition}

\begin{definition}[Guard configuration]
\lean{Turing.MachineFormat.guardCfg}
\label{Turing.MachineFormat.guardCfg}
\leanok
\uses{Turing.Cfg}
Let $b,j\in\mathbb{N}$, let $x$ be a finite Boolean string with $j\leq |x|$, let
$q\in\mathrm{Option}(\mathrm{Fin}(b+1))$, and let $out$ be a finite Boolean string. The
configuration has state $q$ (stored as an active state), input-head position $j+1$, no
work tapes, and output $out$.
\end{definition}

\begin{lemma}[One interior guard step]
\lean{Turing.MachineFormat.guardBlock_step}
\label{Turing.MachineFormat.guardBlock_step}
\leanok
\uses{Turing.Cfg.ext_zero_tapes, Turing.Cfg.inputSymbol, Turing.MachineFormat.guardBlockTM, Turing.MachineFormat.guardCfg, Turing.MachineFormat.guardPhase, Turing.MultiTapeTM.step, Turing.inputMoveAction, Turing.inputSymbolInner, Turing.moveInputPos, Turing.moveInputPos_pos_of_ne_right, Turing.step_of_inputMoveAction}
\difficulty{5}
Let $b$ be a natural number, let $x$ be a finite Boolean string, and let $u,j\in\mathbb
N$ with $u<b+1$ and $j<|x|$. For any Boolean output string $out$, one step of the guard
machine with parameter $b$, from position $j$ in phase $u$, gives the configuration at
position $j+1$ in phase $u+1$, with output unchanged.
\end{lemma}

\begin{lemma}[Guard block run advances by the requested length]
\lean{Turing.MachineFormat.guardBlock_run}
\label{Turing.MachineFormat.guardBlock_run}
\leanok
\uses{Turing.MachineFormat.guardBlockTM, Turing.MachineFormat.guardBlock_step, Turing.MachineFormat.guardCfg, Turing.MachineFormat.guardPhase, Turing.MultiTapeTM.runFrom, Turing.MultiTapeTM.runFrom_succ_eq_step'}
\difficulty{4}
Let $b$ be a natural number, let $x$ be a finite Boolean string, and let $j,u\in\mathbb
N$ with $j\le |x|$, $u\le b+1$, and $j+u\le |x|$. For any Boolean output string $out$,
after $u$ steps of the guard machine with parameter $b$ from position $j$ in phase zero,
the configuration is at position $j+u$ in phase $u$ modulo $b+1$, with output unchanged.
\end{lemma}

\begin{lemma}[Blank input halts the guard scan]
\lean{Turing.MachineFormat.guardBlock_blank}
\label{Turing.MachineFormat.guardBlock_blank}
\leanok
\uses{Turing.Cfg.inputSymbol, Turing.MachineFormat.guardBlockTM, Turing.MachineFormat.guardCfg, Turing.MachineFormat.guardPhase, Turing.MultiTapeTM.step, Turing.moveInputPos_zero, Turing.step_of_inputMoveAction}
\difficulty{3}
For any natural number $b$, Boolean input string $x$, natural number $u$, and Boolean
output string $out$, one step of the guard machine with parameter $b$, started in phase
$u$ at the right blank cell of $x$ with output $out$, yields the configuration in the
halting state at that same cell, with output unchanged.
\end{lemma}

\begin{lemma}[Guard block halts at the input boundary]
\lean{Turing.MachineFormat.guardBlock_short}
\label{Turing.MachineFormat.guardBlock_short}
\leanok
\uses{Turing.MachineFormat.guardBlockTM, Turing.MachineFormat.guardBlock_blank, Turing.MachineFormat.guardBlock_run, Turing.MachineFormat.guardCfg, Turing.MachineFormat.guardPhase, Turing.MultiTapeTM.runFrom, Turing.MultiTapeTM.runFrom_succ_eq_step'}
\difficulty{4}
Let $b$ be a natural number, let $x$ be a finite Boolean string, and let $j\in\mathbb N$
satisfy $j\le |x|$ and $|x|-j<b+1$. For any Boolean output string $out$, after $|x|-j+1$
steps of the guard machine with parameter $b$, started at position $j$ in phase zero,
the configuration is halted at position $|x|$, with output unchanged.
\end{lemma}

\begin{definition}[Guard time bound]
\lean{Turing.MachineFormat.guardTime}
\label{Turing.MachineFormat.guardTime}
\leanok
\uses{Turing.Cfg}
Let $b\in\mathbb{N}$, let $x$ be a finite input string, and let $c$ be a configuration
with no work tapes, Boolean tape alphabet, and states in
$\mathrm{Option}(\mathrm{Fin}(b+1))$. Its guard time is the smaller of $b+1$ and
$|x|+2-c_{\mathrm{in}}$, where $c_{\mathrm{in}}$ is the numerical input-head position.
\end{definition}

\begin{lemma}[Interior guard phases]
\lean{Turing.MachineFormat.guardBlock_interior}
\label{Turing.MachineFormat.guardBlock_interior}
\leanok
\uses{Turing.MachineFormat.guardBlockTM, Turing.MachineFormat.guardBlock_run, Turing.MachineFormat.guardCfg, Turing.MachineFormat.guardPhase, Turing.MachineFormat.guardTime, Turing.MultiTapeTM.runFrom}
\difficulty{5}
Let $b$ be a natural number, let $x$ be a finite Boolean string, and let $j$ be a
natural number with $j\le |x|$. For any Boolean output string $out$ and natural number
$u$ satisfying $0<u<\min(b+1,|x|+2-(j+1))$, after $u$ steps of the guard machine with
parameter $b$ from the configuration at position $j$ in phase zero, the state is some
phase $s$ that is neither phase zero nor the halting state.
\end{lemma}

\begin{definition}[Binary guard machine]
\lean{Turing.MachineFormat.binaryGuardTM}
\label{Turing.MachineFormat.binaryGuardTM}
\leanok
\uses{Turing.CounterLoopState, Turing.MachineFormat.guardBlockTM, Turing.MachineFormat.guardPhase, Turing.MultiTapeTM, Turing.counterLoopTM}
For each $b\in\mathbb{N}$, the binary guard machine is the one-work-tape Boolean machine
whose state set consists of loop states built from $\mathrm{Option}(\mathrm{Fin}(b+1))$.
It is obtained by placing the guard block for $b$ inside the counter-loop machine, with
initial guard phase $0$ and initial counter value $\mathrm{some}(\mathrm{none})$.
\end{definition}

\begin{definition}[Binary guard starting configuration]
\lean{Turing.MachineFormat.binaryGuardStart}
\label{Turing.MachineFormat.binaryGuardStart}
\leanok
\uses{Turing.Cfg, Turing.CounterLoopState, Turing.FinTM.bufferTape, Turing.MachineFormat.guardCfg, Turing.MachineFormat.guardPhase, Turing.counterLoopCfg}
Let $b,j\in\mathbb{N}$, let $x$ be a finite Boolean string with $j\leq |x|$, and let
$out,w$ be finite Boolean strings. The starting configuration has one work tape, Boolean
tape alphabet, and loop states built from $\mathrm{Option}(\mathrm{Fin}(b+1))$; its
state is the decrement phase, its guard component has phase $0$, input-head position
$j+1$, and output $out$, and its work tape is initialized with buffer $w$ at head
position $0$.
\end{definition}

\begin{lemma}[Guard block orbit advances by full blocks]
\lean{Turing.MachineFormat.guardBlock_orbit}
\label{Turing.MachineFormat.guardBlock_orbit}
\leanok
\uses{Turing.MachineFormat.guardBlockTM, Turing.MachineFormat.guardBlock_run, Turing.MachineFormat.guardCfg, Turing.MachineFormat.guardPhase, Turing.MachineFormat.guardTime, Turing.MultiTapeTM.runFrom, Turing.counterLoopOrbit}
\difficulty{5}
Let $b$ be a natural number, let $x$ be a finite Boolean string, and let $j,r\in\mathbb
N$ with $j\le |x|$ and $j+(b+1)r\le |x|$. For any Boolean output string $out$, iterating
the guard machine's counter-loop orbit $r$ times from position $j$ in phase zero gives
the configuration in phase zero at position $j+(b+1)r$, with output unchanged.
\end{lemma}

\begin{lemma}[Guard block segment phases]
\lean{Turing.MachineFormat.guardBlock_segment}
\label{Turing.MachineFormat.guardBlock_segment}
\leanok
\uses{Turing.MachineFormat.guardBlockTM, Turing.MachineFormat.guardBlock_interior, Turing.MachineFormat.guardCfg, Turing.MachineFormat.guardPhase, Turing.MachineFormat.guardTime, Turing.MultiTapeTM.runFrom}
\difficulty{3}
Let $b$ be a natural number, let $x$ be a finite Boolean string, and let $j\in\mathbb N$
with $j\le |x|$. For any Boolean output string $out$, let $c$ be the guard configuration
at position $j$ in phase zero, with output $out$. Its state is phase zero and its guard
time is positive; moreover, for every natural number $u$ strictly between zero and that
guard time, after $u$ steps the machine is in a state that is neither phase zero nor the
halting state.
\end{lemma}

\begin{lemma}[Halting of the binary guard machine]
\lean{Turing.MachineFormat.binaryGuard_done}
\label{Turing.MachineFormat.binaryGuard_done}
\leanok
\uses{Turing.FinTM.bufferTape, Turing.MachineFormat.binaryGuardStart, Turing.MachineFormat.binaryGuardTM, Turing.MachineFormat.guardBlockTM, Turing.MachineFormat.guardBlock_orbit, Turing.MachineFormat.guardBlock_segment, Turing.MachineFormat.guardCfg, Turing.MachineFormat.guardPhase, Turing.MachineFormat.guardTime, Turing.MultiTapeTM.runFrom, Turing.counterLoopOrbit, Turing.counterLoopTM_run_done, Turing.counterLoop_time_le, Turing.counterOverhead}
\difficulty{6}
Let $b,j\in\mathbb{N}$, let $x,w,out$ be finite binary words, and suppose $j\leq |x|$
and $j+(b+1)\operatorname{val}_2(w)\leq |x|$, where $\operatorname{val}_2(w)$ is the
natural number represented by $w$ (with its bits interpreted as $0$ and $1$). Starting
the binary guard machine with parameter $b$ on input $x$, at input position $j$, with
output $out$ and its counter tape initialized with $w$, there is a time $T\leq
5|x|+2|w|+2$ at which the machine is in its done state, its output is still $out$, and
its input-head position is $j+(b+1)\operatorname{val}_2(w)+1$. At every time strictly
before $T$, its state is neither done nor escape.
\end{lemma}

\begin{lemma}[Binary guard escape]
\lean{Turing.MachineFormat.binaryGuard_escape}
\label{Turing.MachineFormat.binaryGuard_escape}
\leanok
\uses{Turing.FinTM.bufferTape, Turing.MachineFormat.binaryGuardStart, Turing.MachineFormat.binaryGuardTM, Turing.MachineFormat.guardBlockTM, Turing.MachineFormat.guardBlock_orbit, Turing.MachineFormat.guardBlock_segment, Turing.MachineFormat.guardBlock_short, Turing.MachineFormat.guardCfg, Turing.MachineFormat.guardPhase, Turing.MachineFormat.guardTime, Turing.MultiTapeTM.runFrom, Turing.counterLoopOrbit, Turing.counterLoopTM_run_escape, Turing.counterOverhead, Turing.counterOverhead_le_of_le}
\difficulty{6}
Let $b,j$ be nonnegative integers, let $x$ be a finite Boolean word with $j\le |x|$, and
let $out$ and $w$ be finite Boolean words. If $|x|<j+(b+1)N$, where $N$ is the natural
number whose binary digits are the entries of $w$, then there is a time $T\le
(b+5)(|x|+1)+2|w|$ at which the binary guard machine, started on $x$ at input-head
position $j$ with output $out$ and counter tape storing $w$, is in the escape state, its
output is still $out$, and its input-head position is $|x|+1$. At every earlier time
$u<T$, its state is neither the done state nor the escape state.
\end{lemma}

\begin{definition}[Canonical Boolean-list check]
\lean{Turing.MachineFormat.canonBitsOK}
\label{Turing.MachineFormat.canonBitsOK}
\leanok
For a finite Boolean string $w$, the check returns its last bit if the string is
nonempty, and returns true if it is empty.
\end{definition}

\begin{lemma}[Boolean fold and final symbol]
\lean{Turing.MachineFormat.canonBits_fold}
\label{Turing.MachineFormat.canonBits_fold}
\leanok
\difficulty{4}
For every Boolean word $w$ and Boolean value $a$, folding the symbols of $w$ from left
to right with the operation that replaces the accumulator by the current symbol,
starting from $a$, yields the last symbol of $w$, or $a$ when $w$ is empty.
\end{lemma}

\begin{lemma}[Canonical bits from the final symbol]
\lean{Turing.MachineFormat.canonBits_of_last}
\label{Turing.MachineFormat.canonBits_of_last}
\leanok
\uses{Turing.MachineFormat.canonBitsOK}
\difficulty{5}
For every Boolean word $w$, if its last symbol, interpreted as true when $w$ is empty,
is true, then the binary expansion of the natural number whose binary digits are the
symbols of $w$ read from right to left is exactly $w$.
\end{lemma}

\begin{lemma}[Canonical expansion has a true final bit]
\lean{Turing.MachineFormat.canonBits_bits}
\label{Turing.MachineFormat.canonBits_bits}
\leanok
\uses{Turing.MachineFormat.canonBitsOK, Turing.codeBitsNat_bits}
\difficulty{5}
For every natural number $n$, the last bit of the binary expansion of $n$, interpreted
as true when that expansion is empty, is true.
\end{lemma}

\begin{lemma}[Characterization of canonical Boolean words]
\lean{Turing.MachineFormat.canonBitsOK_iff}
\label{Turing.MachineFormat.canonBitsOK_iff}
\leanok
\uses{Turing.MachineFormat.canonBitsOK, Turing.MachineFormat.canonBits_bits, Turing.MachineFormat.canonBits_of_last}
\difficulty{3}
For every Boolean word $w$, its last symbol, taken to be true when $w$ is empty, is true
if and only if $w$ is the binary expansion of the natural number obtained by folding its
symbols from right to left as binary digits, with initial value zero.
\end{lemma}

\begin{definition}[Canonical count string]
\lean{Turing.MachineFormat.canonCount}
\label{Turing.MachineFormat.canonCount}
\leanok
\uses{Turing.incFixed}
For a finite Boolean string $w$, append a false bit and apply the fixed-width increment
operation; the resulting canonical count string is that operation's output, or the empty
string if it has no output.
\end{definition}

\begin{lemma}[Incremented Canonical Count]
\lean{Turing.MachineFormat.canonCount_spec}
\label{Turing.MachineFormat.canonCount_spec}
\leanok
\uses{Turing.MachineFormat.canonCount, Turing.incFixed}
\difficulty{4}
For every Boolean word $w$, incrementing the word obtained by appending a zero succeeds,
returns the canonical count word, and increases both its length and its little-endian
binary value by one. That is, the result has length $|w|+1$ and value equal to the value
of $w$ plus one.
\end{lemma}

\begin{definition}[Canonical Boolean-scanning machine]
\lean{Turing.MachineFormat.canonBitsTM}
\label{Turing.MachineFormat.canonBitsTM}
\leanok
\uses{Turing.MultiTapeTM, Turing.ScanPhase, Turing.scanFold, Turing.scanKeep, Turing.scanTM}
For each $k\in\mathbb{N}$ and tape index $i\in\mathrm{Fin}(k)$, this is a $k$-work-tape
Boolean machine with scan states recording a Boolean value. It scans while folding the
input bits by replacing the recorded value with each bit read, keeps the work tapes, and
starts with recorded value true.
\end{definition}

\begin{lemma}[Run of the canonical-bit scan machine]
\lean{Turing.MachineFormat.canonBitsTM_run}
\label{Turing.MachineFormat.canonBitsTM_run}
\leanok
\uses{Turing.Cfg.ofWords, Turing.MachineFormat.canonBitsOK, Turing.MachineFormat.canonBitsTM, Turing.MachineFormat.canonBits_fold, Turing.MultiTapeTM.runFrom, Turing.scanFold, Turing.scanKeep, Turing.scanResult_fold, Turing.scanTM_run_ofCfg}
\difficulty{4}
Let $k\in\mathbb{N}$, let $i<k$ be a work-tape index, let $\alpha$ be a Boolean input
word, and let $w$ assign a Boolean word to each of the $k$ work tapes. Starting from the
configuration with state $(\mathrm{fwd},\mathrm{true})$, input head at position $1$,
those tape contents, all work-tape heads at position $0$, and empty output, the
canonical-bit scan machine reaches after exactly $2|w_i|+2$ steps the configuration with
state $(\mathrm{done},c)$, the same tape contents and otherwise the standard initial
configuration data, where $c$ is true precisely when the last symbol of $w_i$ (or true
if $w_i$ is empty) is true. At every earlier time $t<2|w_i|+2$, its state is not
$(\mathrm{done},b)$ for any Boolean value $b$.
\end{lemma}

\begin{lemma}[Canonical Header Seam Machine]
\lean{Turing.MachineFormat.exists_canonHeaderSeam}
\label{Turing.MachineFormat.exists_canonHeaderSeam}
\leanok
\uses{Turing.Cfg.mapState, Turing.Cfg.ofWords, Turing.FinTM, Turing.MachineFormat.canonHeader, Turing.MachineFormat.exists_canonHeaderCall, Turing.MultiTapeTM.runFrom, Turing.seamReleaseTM, Turing.seamReleaseTM_firstReturn, Turing.stateWord}
\difficulty{5}
For every machine format $F$, there exist a finite Turing machine $C$ over the Boolean
alphabet with at least one work tape, entry and exit states, and a constant
$a\in\mathbb N$. On every Boolean input $\alpha$, it reaches within $t\le
a(|\alpha|+1)$ steps, where $t>0$, the configuration with the exit state, the canonical
header of $\alpha$ on tape zero, and empty words on all other tapes; the input head is
at position $1$ and the output is empty. At every earlier time $u<t$ its state is not
the exit state.
\end{lemma}

\begin{lemma}[Linear-time computation of the canonical count]
\lean{Turing.MachineFormat.exists_canonCountTM}
\label{Turing.MachineFormat.exists_canonCountTM}
\leanok
\uses{Turing.FinTM, Turing.FinTM.ComputesFunInTime, Turing.FinTM.ComputesInTime.mono, Turing.FinTM.computesFunInTime_appendBit, Turing.FinTM.computesFunInTime_comp, Turing.FinTM.computesFunInTime_incFixed, Turing.MachineFormat.canonCount}
\difficulty{3}
For every machine format $F$, there exist a finite Turing machine $M$ over the Boolean
alphabet and a natural number $a$ such that, on every Boolean input word $w$, $M$
computes the canonical count of $w$ within $a(|w|+1)$ steps.
\end{lemma}

\begin{definition}[Canonical verdict machine]
\lean{Turing.MachineFormat.canonVerdictTM}
\label{Turing.MachineFormat.canonVerdictTM}
\leanok
\uses{Turing.MachineFormat.canonBitsTM, Turing.MultiTapeTM, Turing.ScanPhase}
For each $k\in\mathbb{N}$ and tape index $i\in\mathrm{Fin}(k)$, this is a $k$-work-tape
Boolean machine with scan states recording a Boolean value, initially true in the
forward phase. In a completed state carrying verdict
$b\in\{\mathrm{false},\mathrm{true}\}$, it stays in place, leaves every work tape
unchanged, outputs $b$, and halts; in all other states it follows the canonical
Boolean-scanning machine's transition function.
\end{definition}

\begin{lemma}[Run of the canonical-bit verdict machine]
\lean{Turing.MachineFormat.canonVerdictTM_run}
\label{Turing.MachineFormat.canonVerdictTM_run}
\leanok
\uses{Turing.Action.apply, Turing.Action.mapState, Turing.Cfg.mapState, Turing.Cfg.ofWords, Turing.MachineFormat.canonBitsOK, Turing.MachineFormat.canonBitsTM, Turing.MachineFormat.canonBitsTM_run, Turing.MachineFormat.canonVerdictTM, Turing.MultiTapeTM.runFrom, Turing.MultiTapeTM.runFrom_mapState_of_agreeOn, Turing.MultiTapeTM.runFrom_succ_eq_step', Turing.MultiTapeTM.step, Turing.ScanPhase}
\difficulty{5}
Let $k$ be a natural number, let $i$ be one of $k$ work tapes, let $\alpha$ be a Boolean
input word, and let $w$ assign a Boolean word to each work tape. Starting in the
forward-pass state with the words $w$ on the tapes, after exactly $2|w(i)|+3$ steps the
verdict machine has halted and its output is the single bit equal to the last bit of
$w(i)$, or $1$ if $w(i)$ is empty.
\end{lemma}

\begin{lemma}[Linear-Time Canonical Header Test]
\lean{Turing.MachineFormat.exists_canonHeaderBitsTM}
\label{Turing.MachineFormat.exists_canonHeaderBitsTM}
\leanok
\uses{Turing.Cfg.init, Turing.Cfg.mapState, Turing.Cfg.ofWords, Turing.FinTM, Turing.FinTM.ComputesFunInTime, Turing.FinTM.ComputesInTime, Turing.FinTM.ComputesInTime.mono, Turing.FinTM.bufferTape_nil, Turing.FinTM.computesInTime_iff, Turing.MachineFormat.canonBitsOK, Turing.MachineFormat.canonHeader, Turing.MachineFormat.canonHeader_length, Turing.MachineFormat.canonVerdictTM, Turing.MachineFormat.canonVerdictTM_run, Turing.MachineFormat.exists_canonHeaderSeam, Turing.MultiTapeTM.initCfg, Turing.MultiTapeTM.runFrom, Turing.ScanPhase, Turing.seamCompTM, Turing.seamCompTM_run_ofCfg, Turing.stateWord}
\difficulty{5}
For every machine format $F$, there exist a finite Turing machine $M$ over the Boolean
alphabet and a constant $a\in\mathbb N$ such that, on every Boolean word $\alpha$,
$M$ halts after exactly $a(|\alpha|+1)$ steps and outputs the one-bit word whose bit is
true exactly when the last bit of the parsed header (defaulting to true if that header
is empty) is true.
\end{lemma}

\begin{definition}[Canonical header validity]
\lean{Turing.MachineFormat.canonHeaderOK}
\label{Turing.MachineFormat.canonHeaderOK}
\leanok
\uses{Turing.MachineFormat.canonBitsOK, Turing.MachineFormat.canonHeader, Turing.pairDecode}
A bit string is canonically valid as a header exactly when it decodes as a pair and the
canonical header extracted from it passes the bit-validity check.
\end{definition}

\begin{lemma}[Characterization of valid canonical headers]
\lean{Turing.MachineFormat.canonHeaderOK_iff}
\label{Turing.MachineFormat.canonHeaderOK_iff}
\leanok
\uses{Turing.MachineFormat.canonBitsOK_iff, Turing.MachineFormat.canonHeader, Turing.MachineFormat.canonHeaderOK, Turing.codeBitsNat, Turing.pairDecode}
\difficulty{4}
For every Boolean word $\alpha$, the canonical-header validity test returns true if and
only if $\alpha$ decodes as a pair $(bits,rest)$ and $bits$ is the canonical binary
expansion of the natural number denoted by $bits$.
\end{lemma}

\begin{lemma}[Malformed header rejection]
\lean{Turing.MachineFormat.scan_eq_none_of_canonHeaderOK}
\label{Turing.MachineFormat.scan_eq_none_of_canonHeaderOK}
\leanok
\uses{Turing.CodeFormat.scan, Turing.MachineFormat, Turing.MachineFormat.canonHeaderOK, Turing.MachineFormat.canonHeaderOK_iff, Turing.codeBitsNat, Turing.pairDecode}
\difficulty{4}
Let $F$ be a machine format and let $\alpha$ be a finite bit string. If either $\alpha$
has no aligned-pair parse, or its parsed header does not end in $1$ (with the empty
header treated as passing this last-bit check), then scanning $\alpha$ fails.
\end{lemma}

\begin{lemma}[Linear-Time Canonical Header Validity Test]
\lean{Turing.MachineFormat.exists_canonHeaderOKTM}
\label{Turing.MachineFormat.exists_canonHeaderOKTM}
\leanok
\uses{Turing.FinTM, Turing.FinTM.ComputesFunInTime, Turing.FinTM.ComputesInTime, Turing.FinTM.ComputesInTime.mono, Turing.FinTM.computesFunInTime_cond, Turing.FinTM.computesFunInTime_const, Turing.FinTM.computesFunInTime_pairValid, Turing.MachineFormat.canonBitsOK, Turing.MachineFormat.canonHeader, Turing.MachineFormat.canonHeaderOK, Turing.MachineFormat.exists_canonHeaderBitsTM, Turing.pairDecode}
\difficulty{5}
For every machine format $F$, there exist a finite Turing machine $M$ over the Boolean
alphabet and a constant $a\in\mathbb N$ such that, on every Boolean input $\alpha$,
$M$ halts after exactly $a(|\alpha|+1)$ steps and outputs the one-bit word indicating
whether $\alpha$ has a valid pair parse and its parsed header is canonical, meaning its
last bit is true (with the empty header accepted).
\end{lemma}

\begin{definition}[Canonical guard block size]
\lean{Turing.MachineFormat.canonGuardBlock}
\label{Turing.MachineFormat.canonGuardBlock}
\leanok
\uses{Turing.MachineFormat}
For a machine format with guard constant $g$ and $k$ work tapes, the canonical guard
block size is $\max\bigl(g,(4k+5)3^{k+1}\bigr)$.
\end{definition}

\begin{lemma}[Positivity of the Guard Block]
\lean{Turing.MachineFormat.canonGuardBlock_pos}
\label{Turing.MachineFormat.canonGuardBlock_pos}
\leanok
\uses{Turing.MachineFormat, Turing.MachineFormat.canonGuardBlock}
\difficulty{3}
For every machine format $F$, its guard-block bound is positive.
\end{lemma}

\begin{lemma}[Short Codes Decode to the Fallback]
\lean{Turing.MachineFormat.decode_of_canonGuard}
\label{Turing.MachineFormat.decode_of_canonGuard}
\leanok
\uses{Turing.CodeFormat.decode, Turing.CodeFormat.decode_of_guard_lt, Turing.MachineFormat, Turing.MachineFormat.canonGuardBlock, Turing.MachineFormat.decode_of_grammar_guard, Turing.codeBitsNat, Turing.pairDecode}
\difficulty{4}
Let $F$ be a machine format, and let $\alpha$, $b$, and $r$ be finite bit strings such
that parsing $\alpha$ as a paired encoding returns the pair $(b,r)$. Let $m=\sum_{j<|b|}
b_j2^j$ be the natural number represented by $b$, with its bits read least-significant
first, and let $B=\max\bigl(g,(4k+5)3^{k+1}\bigr)$, where $g$ is the format's guard and
$k$ is its number of work tapes. If $|r|<B(m+1)$, then decoding $\alpha$ in the format
$F$ gives its fallback code.
\end{lemma}

\begin{lemma}[Bounded run of the canonicality guard]
\lean{Turing.MachineFormat.canonGuard_run}
\label{Turing.MachineFormat.canonGuard_run}
\leanok
\uses{Turing.MachineFormat, Turing.MachineFormat.binaryGuardStart, Turing.MachineFormat.binaryGuardTM, Turing.MachineFormat.binaryGuard_done, Turing.MachineFormat.binaryGuard_escape, Turing.MachineFormat.canonCount, Turing.MachineFormat.canonCount_spec, Turing.MachineFormat.canonGuardBlock, Turing.MachineFormat.canonGuardBlock_pos, Turing.MultiTapeTM.runFrom, Turing.codeBitsNat, Turing.eq_pairEncode_of_pairDecode, Turing.foldr_bit_eq_ofDigits, Turing.length_pairEncode, Turing.pairDecode}
\difficulty{6}
Let $F$ be a machine format, and let $\alpha$, $bits$, and $rest$ be Boolean words such
that parsing $\alpha$ as a paired encoding returns $(bits,rest)$. Set
$B=\max\{F.\mathrm{guard},(4F.\mathrm{k}+5)3^{F.\mathrm{k}+1}\}$ and $j=2|bits|+2$.
There exist a position bound $j\le|\alpha|$ and a time
$T\le(B+4)(|\alpha|+1)+2(|bits|+1)$ such that the binary guard machine with block
parameter $B-1$, started at position $j$ with empty output and counter tape containing
the fixed-width increment of $bits$ followed by a zero, is in its done state at time $T$
if $B(\operatorname{value}(bits)+1)\le|rest|$, and in its escape state otherwise; its
output is empty, and its input-head position is respectively
$j+B(\operatorname{value}(bits)+1)+1$ or $|\alpha|+1$. Neither terminal state occurs at
any time strictly before $T$.
\end{lemma}

\begin{lemma}[Strict installation of a computed function]
\lean{Turing.MachineFormat.canonStrictInstall}
\label{Turing.MachineFormat.canonStrictInstall}
\leanok
\uses{Turing.Cfg.mapState, Turing.Cfg.ofWords, Turing.FinTM, Turing.FinTM.ComputesFunInTime, Turing.FinTM.exists_installCallTM, Turing.MultiTapeTM.runFrom, Turing.seamReleaseTM, Turing.seamReleaseTM_firstReturn, Turing.stateWord}
\difficulty{5}
Let $M$ be a finite Turing machine over the Boolean alphabet computing a function $f$ in
time $T$. Then there exist a finite Boolean Turing machine $C$, two states of $C$ called
entry and exit, and a natural number $a$, with at least one work tape, such that for
every pair of Boolean words $x$ and $\mathrm{arg}$, the machine started on input $x$ in
the entry state, with $\mathrm{arg}$ on tape $0$ and all other work tapes empty, reaches
after some positive number $t$ of steps at most
$a\bigl(T(|\mathrm{arg}|)+|\mathrm{arg}|+|f(\mathrm{arg})|+1\bigr)$ the configuration
with state exit, $f(\mathrm{arg})$ on tape $0$, and all other work tapes empty. At every
earlier time $u<t$, its state is different from exit.
\end{lemma}

\begin{lemma}[Embedded Configuration Equality]
\lean{Turing.MachineFormat.canonEmbedEq}
\label{Turing.MachineFormat.canonEmbedEq}
\leanok
\uses{Turing.Cfg, Turing.MultiTapeTM, Turing.embedEmitCfg, Turing.embedEmitCfg_selected_pos, Turing.embedEmitCfg_selected_tape, Turing.embedEmitTM_frame}
\difficulty{4}
Let $m,k$ be natural numbers, let $S$ be a type, and let $x$ be a Boolean word. Let $M$
be an $m$-work-tape Turing machine, let $\iota$ inject the $m$ work-tape indices into
the $k$ indices, and let $c,d$ be configurations on input $x$ with respectively $m$ and
$k$ work tapes. Suppose their states, input-head positions, and outputs agree; for each
$i<m$, the contents and head position of tape $i$ in $c$ agree with those of tape
$\iota(i)$ in $d$; and outside the image of $\iota$, the supplied tape contents and head
positions agree with those in $d$. Then embedding $c$ into the $k$-tape configuration
with the supplied contents and head positions, and with empty prefix output, gives
exactly $d$.
\end{lemma}

\begin{lemma}[Embedded Word Configuration Equality]
\lean{Turing.MachineFormat.canonEmbedWords}
\label{Turing.MachineFormat.canonEmbedWords}
\leanok
\uses{Turing.Cfg.ofWords, Turing.FinTM.bufferTape, Turing.MachineFormat.canonEmbedEq, Turing.MultiTapeTM, Turing.embedEmitCfg}
\difficulty{3}
Let $m,k$ be natural numbers, let $S$ be a type, and let $x$ be a Boolean input word.
Let $M$ be an $m$-work-tape Turing machine over Boolean symbols and states $S$, let
$\iota$ embed the $m$ tape indices into the $k$ tape indices, and let $q$ be a state.
Suppose $v$ assigns a Boolean word to each of the $m$ tapes, and the family
$\mathrm{wanted}$ of $k$ words agrees with $v$ on the embedded indices. Suppose also
that the family $\mathrm{ambient}$ agrees with $\mathrm{wanted}$ outside the image of
$\iota$. Then embedding the configuration formed from $q$ and $v$, using the ambient
words as tape contents and zero head positions on the other tapes, gives the
configuration formed from $q$ and $\mathrm{wanted}$.
\end{lemma}

\begin{lemma}[Installing a computed function at a tape seam]
\lean{Turing.MachineFormat.canonPreserveInstall}
\label{Turing.MachineFormat.canonPreserveInstall}
\leanok
\uses{Turing.Cfg.mapState, Turing.Cfg.ofWords, Turing.FinTM, Turing.FinTM.ComputesFunInTime, Turing.FinTM.bufferTape, Turing.MachineFormat.canonEmbedWords, Turing.MachineFormat.canonStrictInstall, Turing.MultiTapeTM.runFrom, Turing.SweepPhase, Turing.copyTM, Turing.copyTM_run, Turing.embedEmitCfg, Turing.embedEmitTM, Turing.embedEmitTM_runFrom, Turing.seamCompTM, Turing.seamCompTM_firstReturn_ofCfg, Turing.seamCompTM_run_ofCfg, Turing.stateWord}
\difficulty{6}
Let $F$ be a machine format, let $M$ be a finite Turing machine over Boolean symbols,
let $f$ map Boolean words to Boolean words, and let $T$ map natural numbers to natural
numbers. Suppose that on every input word $arg$, $M$ halts within $T(|arg|)$ steps with
output $f(arg)$. Then there exist a finite Boolean Turing machine $C$, a destination
work tape $dst$, entry and exit states, and a natural number $a$, with $dst$ having
index $1$, such that for every input word $x$ and argument word $arg$, starting $C$ in
its entry state with $arg$ on tape $0$ and all other tapes empty, it reaches within
$a(T(|arg|)+|arg|+|f(arg)|+1)$ steps a configuration in the exit state whose tape words
are the initial assignment with $f(arg)$ installed on tape $dst$. The exit state is not
reached at any earlier time.
\end{lemma}

\begin{lemma}[Canonical Header Count Seam]
\lean{Turing.MachineFormat.exists_canonHeaderCountSeam}
\label{Turing.MachineFormat.exists_canonHeaderCountSeam}
\leanok
\uses{Turing.Cfg.mapState, Turing.Cfg.ofWords, Turing.FinTM, Turing.FinTM.bufferTape, Turing.MachineFormat.canonCount, Turing.MachineFormat.canonCount_spec, Turing.MachineFormat.canonEmbedWords, Turing.MachineFormat.canonHeader, Turing.MachineFormat.canonHeader_length, Turing.MachineFormat.canonPreserveInstall, Turing.MachineFormat.exists_canonCountTM, Turing.MachineFormat.exists_canonHeaderSeam, Turing.MultiTapeTM.runFrom, Turing.embedEmitCfg, Turing.embedEmitTM, Turing.embedEmitTM_runFrom, Turing.seamCompTM, Turing.seamCompTM_firstReturn_ofCfg, Turing.seamCompTM_run_ofCfg, Turing.stateWord}
\difficulty{6}
For every machine format $F$, there exist a finite Turing machine $C$ over the Boolean
alphabet, a work-tape index $d$ whose index is $1$, entry and exit states, and a
constant $a\in\mathbb N$. For every Boolean input $\alpha$, at some time $t\le
a(|\alpha|+1)$ the machine reaches the exit state with the canonical header of
$\alpha$ on tape zero, its fixed-width increment (with overflow replaced by the empty
word) on tape $d$, and empty words on all other tapes; at every earlier time $u<t$ its
state is not the exit state.
\end{lemma}

\begin{definition}[Canonical offset code]
\lean{Turing.MachineFormat.canonOffset}
\label{Turing.MachineFormat.canonOffset}
\leanok
\uses{Turing.MachineFormat.canonCount}
For a bit string $w$, its canonical offset code is a leading $1$ followed by the binary
representation of the length of $w$.
\end{definition}

\begin{lemma}[Value and length of the canonical offset]
\lean{Turing.MachineFormat.canonOffset_spec}
\label{Turing.MachineFormat.canonOffset_spec}
\leanok
\uses{Turing.MachineFormat.canonCount, Turing.MachineFormat.canonCount_spec, Turing.MachineFormat.canonOffset, Turing.codeBitsNat, Turing.codeBitsNat_bits, Turing.foldr_bit_eq_ofDigits}
\difficulty{3}
For every Boolean word $w$, the canonical offset word has binary digit value $2|w|+3$,
and its length is at most $|w|+2$.
\end{lemma}

\begin{lemma}[Linear-time canonical offset]
\lean{Turing.MachineFormat.exists_canonOffsetTM}
\label{Turing.MachineFormat.exists_canonOffsetTM}
\leanok
\uses{Turing.FinTM, Turing.FinTM.ComputesFunInTime, Turing.FinTM.ComputesInTime.mono, Turing.FinTM.computesFunInTime_comp, Turing.FinTM.computesFunInTime_lengthBits, Turing.FinTM.computesFunInTime_prepend, Turing.MachineFormat.canonCount, Turing.MachineFormat.canonOffset, Turing.MachineFormat.exists_canonCountTM}
\difficulty{4}
There are a finite Turing machine $M$ over the Boolean alphabet and a natural number $a$
such that, on every Boolean word $w$, $M$ halts after $a(|w|+1)$ steps with output equal
to the canonical offset of $w$.
\end{lemma}

\begin{definition}[Canonical ready words]
\lean{Turing.MachineFormat.canonReadyWords}
\label{Turing.MachineFormat.canonReadyWords}
\leanok
\uses{Turing.MachineFormat.canonCount, Turing.MachineFormat.canonOffset}
Given $k$ and a header bit string $h$, define a word for each index $j<k$: it is $h$
when $j=0$, the canonical count code of $h$ when $j=1$, the canonical offset code of $h$
when $j=2$, and the empty string for every other index.
\end{definition}

\begin{lemma}[Canonical ready seam computation]
\lean{Turing.MachineFormat.exists_canonReadySeam}
\label{Turing.MachineFormat.exists_canonReadySeam}
\leanok
\uses{Turing.Cfg.mapState, Turing.Cfg.ofWords, Turing.FinTM, Turing.FinTM.bufferTape, Turing.MachineFormat.canonCount, Turing.MachineFormat.canonEmbedWords, Turing.MachineFormat.canonHeader, Turing.MachineFormat.canonHeader_length, Turing.MachineFormat.canonOffset, Turing.MachineFormat.canonOffset_spec, Turing.MachineFormat.canonPreserveInstall, Turing.MachineFormat.canonReadyWords, Turing.MachineFormat.exists_canonHeaderCountSeam, Turing.MachineFormat.exists_canonOffsetTM, Turing.MultiTapeTM.runFrom, Turing.embedEmitCfg, Turing.embedEmitTM, Turing.embedEmitTM_runFrom, Turing.seamCompTM, Turing.seamCompTM_firstReturn_ofCfg, Turing.seamCompTM_run_ofCfg, Turing.stateWord}
\difficulty{6}
There are a finite Turing machine $C$ over the Boolean alphabet, entry and exit states,
and a natural number $a$, with at least four work tapes, such that for every Boolean
word $\alpha$ there is a time $t\le a(|\alpha|+1)$ at which the run from the
configuration with state equal to the entry state, input $\alpha$, blank work tapes, and
empty output has not reached the exit state at any earlier time and equals the
configuration with state equal to the exit state, input head at position $1$, empty
output, and tape contents given by the canonical ready words for the canonical header of
$\alpha$, with all work-tape heads at position zero.
\end{lemma}

\begin{lemma}[Canonical seek on a tape bank]
\lean{Turing.MachineFormat.exists_canonSeekOnBank}
\label{Turing.MachineFormat.exists_canonSeekOnBank}
\leanok
\uses{Turing.Cfg, Turing.FinTM.bufferTape, Turing.MachineFormat.canonEmbedEq, Turing.MultiTapeTM, Turing.MultiTapeTM.runFrom, Turing.embedEmitCfg, Turing.embedEmitTM, Turing.embedEmitTM_runFrom, Turing.exists_inputSeekTM}
\difficulty{4}
Let $F$ be a machine format, let $k$ be a natural number, and let
$\iota:\{0,1\}\hookrightarrow\{0,\ldots,k-1\}$ be an injection. There exist a finite
state set $S$ with decidable equality, a deterministic Turing machine $R$ with $k$ work
tapes over the Boolean alphabet, entry and exit states in $S$, and a natural number $a$
such that the following holds. For every Boolean input word $x$, natural number $p$,
Boolean word $w$ whose base-two value is $p$, and configuration $d$ on input $x$, assume
$p\le |x|+1$, the state of $d$ is the entry state, tape $\iota(0)$ contains the buffer
representation of $w$ at every integer offset $q$ with $-1\le q\le |w|$, and tape
$\iota(1)$ is blank at every integer position. Then for some $t\le
a(|w|+d.inputPos.val+p+1)$, the run of $R$ from $d$ has not reached the exit state at
any time $u<t$, and at time $t$ it is exactly $d$ with its state changed to the exit
state and its input-head position set to $p$.
\end{lemma}

\begin{lemma}[A machine for the canonical guard check]
\lean{Turing.MachineFormat.exists_canonGuardFront}
\label{Turing.MachineFormat.exists_canonGuardFront}
\leanok
\uses{Turing.Cfg, Turing.Cfg.mapState, Turing.Cfg.ofWords, Turing.CounterLoopState, Turing.FinTM, Turing.FinTM.bufferTape, Turing.MachineFormat, Turing.MachineFormat.binaryGuardStart, Turing.MachineFormat.binaryGuardTM, Turing.MachineFormat.canonCount, Turing.MachineFormat.canonEmbedEq, Turing.MachineFormat.canonGuardBlock, Turing.MachineFormat.canonGuard_run, Turing.MachineFormat.canonHeader, Turing.MachineFormat.canonHeader_length, Turing.MachineFormat.canonOffset, Turing.MachineFormat.canonOffset_spec, Turing.MachineFormat.canonReadyWords, Turing.MachineFormat.exists_canonReadySeam, Turing.MachineFormat.exists_canonSeekOnBank, Turing.MachineFormat.guardCfg, Turing.MultiTapeTM.runFrom, Turing.codeBitsNat, Turing.counter_cfg_eq, Turing.embedEmitCfg, Turing.embedEmitTM, Turing.embedEmitTM_frame, Turing.embedEmitTM_runFrom, Turing.eq_pairEncode_of_pairDecode, Turing.length_pairEncode, Turing.pairDecode, Turing.seamCompTM, Turing.seamCompTM_runFrom_left, Turing.seamCompTM_run_ofCfg}
\difficulty{7}
For every machine format $F$, there exist a finite Boolean-alphabet Turing machine $C$,
three pairwise distinct states called entry, yes, and no, and a natural number $a$, with
at least four work tapes, such that for any Boolean strings $\alpha$, $bits$, and $rest$
for which $\alpha$ parses as $bits$ and $rest$, the machine started on $\alpha$ with its
entry state and blank work tapes reaches, at some time $t\le a(|\alpha|+1)$, the yes
state if $\max\{guard,(4k+5)3^{k+1}\}(v+1)\le |rest|$ and the no state otherwise, where
$k$ is the work-tape count of $F$ and $v$ is the binary value of $bits$. It has produced
no output; its input-head position is $2|bits|+3$ in the yes case and $|\alpha|+1$ in
the no case. Every work tape except tape $1$ is at position zero and contains the
canonical count word for $bits$ on tape $0$, the canonical offset word for $bits$ on
tape $2$, and the empty word on all remaining tapes.
\end{lemma}

\begin{definition}[Canonical Boolean pop machine]
\lean{Turing.MachineFormat.canonBoolPopTM}
\label{Turing.MachineFormat.canonBoolPopTM}
\leanok
\uses{Turing.MultiTapeTM, Turing.inputMoveAction, Turing.tapeAction}
For $k$ work tapes and a selected tape $i<k$, define a machine with states
$\{\mathrm{halt},0,1\}$ whose initial state is halt. In the halt state it stays on the
input tape, writes blank on tape $i$, and outputs the bit currently scanned on tape $i$,
treating a blank as $0$; in state $b\in\{0,1\}$ it stays on the input tape and outputs
$b$.
\end{definition}

\begin{lemma}[One-step Boolean tape pop]
\lean{Turing.MachineFormat.canonBoolPop_run}
\label{Turing.MachineFormat.canonBoolPop_run}
\leanok
\uses{Turing.Action.apply, Turing.Cfg.ofWords, Turing.Cfg.workTapeSymbols, Turing.FinTM.bufferTape, Turing.FinTM.bufferTape_append, Turing.MachineFormat.canonBoolPopTM, Turing.MultiTapeTM.runFrom, Turing.MultiTapeTM.runFrom_succ_eq_step, Turing.MultiTapeTM.runFrom_zero, Turing.MultiTapeTM.step, Turing.stateWord, Turing.tapeAction}
\difficulty{4}
Let $k$ be a positive integer, let $b$ be a Boolean value, and let $x$ be a Boolean
input word. Starting with the machine in its initial state and with $[b]$ on work tape
$0$ and every other work tape empty, after one step its state is $b$ and all work tapes
are empty. The input head and work-tape heads remain at their initial positions, and the
output tape remains empty.
\end{lemma}

\begin{lemma}[A machine branching on a computed Boolean value]
\lean{Turing.MachineFormat.exists_canonBoolBranch}
\label{Turing.MachineFormat.exists_canonBoolBranch}
\leanok
\uses{Turing.Cfg, Turing.Cfg.mapState, Turing.Cfg.ofWords, Turing.FinTM, Turing.FinTM.ComputesFunInTime, Turing.FinTM.bufferTape, Turing.MachineFormat.canonBoolPopTM, Turing.MachineFormat.canonBoolPop_run, Turing.MachineFormat.canonStrictInstall, Turing.MultiTapeTM.runFrom, Turing.SweepPhase, Turing.inputCopyTM, Turing.inputCopyTM_run_ofCfg, Turing.seamCompTM, Turing.seamCompTM_firstReturn_ofCfg, Turing.seamCompTM_runFrom_left, Turing.seamCompTM_run_ofCfg, Turing.stateWord}
\difficulty{6}
Let $F$ be a machine format, let $M$ be a finite Turing machine over the Boolean
alphabet, and let $p$ map Boolean strings to Boolean values. Suppose $M$ computes the
one-symbol output $p(x)$ on every input $x$ within $c(|x|+1)$ steps for some natural
number $c$. Then there are a finite Boolean-alphabet Turing machine $C$, an entry state,
an injective assignment of exit states to Boolean values, and a natural number $a$, with
at least one work tape, such that for every input $x$ there is a time $t\le a(|x|+1)$ at
which $C$, started in its entry state with all work tapes blank, is in the exit state
assigned to $p(x)$ with all work tapes blank; before time $t$, it is never in either
exit state.
\end{lemma}

\begin{lemma}[Canonical Header Branch Machine]
\lean{Turing.MachineFormat.exists_canonHeaderBranch}
\label{Turing.MachineFormat.exists_canonHeaderBranch}
\leanok
\uses{Turing.Cfg.ofWords, Turing.FinTM, Turing.MachineFormat.canonHeaderOK, Turing.MachineFormat.exists_canonBoolBranch, Turing.MachineFormat.exists_canonHeaderOKTM, Turing.MultiTapeTM.runFrom}
\difficulty{1}
For every machine format $F$, there exist a finite Turing machine $C$ over the Boolean
alphabet with at least one work tape, an entry state, an exit state, and a constant
$a\in\mathbb N$. On every Boolean input $x$, it reaches within $t\le a(|x|+1)$ steps,
with $t>0$, the configuration in the exit state whose tape words consist of the
canonical header of $x$ on tape zero and empty words on all other tapes; the input head
is at position $1$ and the output is empty. At every time $u<t$ its state is not the
exit state.
\end{lemma}

\begin{lemma}[Fallback on a noncanonical header]
\lean{Turing.MachineFormat.decode_of_canonHeader}
\label{Turing.MachineFormat.decode_of_canonHeader}
\leanok
\uses{Turing.CodeFormat.decode, Turing.CodeFormat.decode_of_parseFull_none, Turing.CodeFormat.scan_full, Turing.MachineFormat, Turing.MachineFormat.canonHeaderOK, Turing.MachineFormat.scan_eq_none_of_canonHeaderOK}
\difficulty{3}
For every machine format $F$ and Boolean string $x$, if the canonical header extracted
from $x$ is invalid, then decoding $x$ in $F$ returns the format’s fallback code.
\end{lemma}

\begin{lemma}[Avoiding a target state across a machine seam]
\lean{Turing.MachineFormat.canonSeam_avoids}
\label{Turing.MachineFormat.canonSeam_avoids}
\leanok
\uses{Turing.Cfg, Turing.Cfg.mapState, Turing.MultiTapeTM, Turing.MultiTapeTM.runFrom, Turing.seamCompTM, Turing.seamCompTM_runFrom_left, Turing.seamCompTM_run_ofCfg}
\difficulty{4}
Let $A$ and $B$ be deterministic Turing machines with $k$ work tapes over Boolean
symbols, with state sets $S$ and $R$, respectively; choose a state $leave\in S$, a state
$enter\in R$, and a target state in $S\oplus R$. For a Boolean input $x$, configurations
$d,e$ of $A$ on $x$, and times $t,v\in\mathbb N$, suppose $A$ runs from $d$ to $e$ in
$t$ steps, $e$ is in state $leave$, and $leave$ is not visited at any earlier time.
Suppose also that the configurations of this run, relabeled into the left side of
$S\oplus R$, never have the target state through time $t$, and that the run of $B$ from
$e$ with its state replaced by $enter$, relabeled into the right side, avoids the target
state for its first $v$ steps. Then the machine that runs $A$ until $leave$ and
continues with $B$ from $enter$ avoids the target state for every time strictly less
than $t+1+v$, when started from $d$ relabeled into the left side.
\end{lemma}

\begin{definition}[Canonical counter bank embedding]
\lean{Turing.MachineFormat.canonCounterBank}
\label{Turing.MachineFormat.canonCounterBank}
\leanok
For $k$ work tapes and a selected tape $i<k$, the canonical counter bank embeds the
singleton set into the $k$ tape indices by sending its sole element to $i$.
\end{definition}

\begin{lemma}[Bounded guard computation on one tape]
\lean{Turing.MachineFormat.canonGuard_onBank}
\label{Turing.MachineFormat.canonGuard_onBank}
\leanok
\uses{Turing.Cfg, Turing.CounterLoopState, Turing.FinTM.bufferTape, Turing.MachineFormat, Turing.MachineFormat.binaryGuardStart, Turing.MachineFormat.binaryGuardTM, Turing.MachineFormat.canonCount, Turing.MachineFormat.canonCounterBank, Turing.MachineFormat.canonEmbedEq, Turing.MachineFormat.canonGuardBlock, Turing.MachineFormat.canonGuard_run, Turing.MachineFormat.canonHeader, Turing.MachineFormat.canonHeader_length, Turing.MachineFormat.guardCfg, Turing.MultiTapeTM.runFrom, Turing.codeBitsNat, Turing.counter_cfg_eq, Turing.embedEmitCfg, Turing.embedEmitTM, Turing.embedEmitTM_frame, Turing.embedEmitTM_runFrom, Turing.pairDecode}
\difficulty{5}
Let $k\in\mathbb N$, let $i$ be one of the $k$ work tapes, and let $F$ be a machine
format. Let $x$, $bits$, and $rest$ be Boolean words such that paired decoding of $x$
yields $(bits,rest)$. Set $B=\max\{F.guard,(4F.k+5)3^{F.k+1}\}$, and let $d$ be a
configuration on input $x$ whose state is the running decrement phase, whose input head
is at $2|bits|+3$, whose tape $i$ stores the word obtained by applying fixed-width
increment to $bits$ with a zero appended (using the empty word if increment overflows),
whose head on tape $i$ is at $0$, and whose output is empty. When the guard counter
machine with block size $B$ is embedded on tape $i$ and run from $d$, there is a time
$t\le (B+6)(|x|+1)$ at which its state is done if $B(\operatorname{value}(bits)+1)\le
|rest|$, and escape otherwise; its output is empty and its input head is respectively at
$2|bits|+3+B(\operatorname{value}(bits)+1)$ or $|x|+1$. At every earlier time $u<t$ its
state is neither done nor escape, and every tape other than $i$ and its head position
are unchanged from $d$.
\end{lemma}

\begin{lemma}[Guard-Checking Machine from Ready Tapes]
\lean{Turing.MachineFormat.exists_canonGuardReady}
\label{Turing.MachineFormat.exists_canonGuardReady}
\leanok
\uses{Turing.Cfg, Turing.Cfg.mapState, Turing.Cfg.ofWords, Turing.CounterLoopState, Turing.FinTM.bufferTape, Turing.MachineFormat, Turing.MachineFormat.binaryGuardTM, Turing.MachineFormat.canonCounterBank, Turing.MachineFormat.canonGuardBlock, Turing.MachineFormat.canonGuard_onBank, Turing.MachineFormat.canonHeader, Turing.MachineFormat.canonHeader_length, Turing.MachineFormat.canonOffset, Turing.MachineFormat.canonOffset_spec, Turing.MachineFormat.canonReadyWords, Turing.MachineFormat.canonSeam_avoids, Turing.MachineFormat.exists_canonSeekOnBank, Turing.MultiTapeTM, Turing.MultiTapeTM.runFrom, Turing.codeBitsNat, Turing.embedEmitTM, Turing.eq_pairEncode_of_pairDecode, Turing.length_pairEncode, Turing.pairDecode, Turing.seamCompTM, Turing.seamCompTM_runFrom_left, Turing.seamCompTM_run_ofCfg}
\difficulty{7}
For every machine format $F$, every $k\ge4$, there exist a finite state set $S$, a
deterministic Turing machine with $k$ work tapes and states $S$, an entry state, an
injective map from the Boolean values to states, and a constant $a\in\mathbb N$. For
all Boolean words $x,bits,rest$ parsed from $x$ as $(bits,rest)$, the machine, started
on $x$ with the prescribed ready words on its tapes, reaches within $t\le a(|x|+1)$
steps the exit state indexed by whether $B_F(\operatorname{val}(bits)+1)\le |rest|$,
where $B_F=\max(\operatorname{guard}(F),(4k+5)3^{k+1})$; at that time its output is
empty, its input-head position is $2|bits|+3$ when the inequality holds and $|x|+1$
otherwise, and every tape except tape $1$ retains its prescribed ready word with its
head at zero. No earlier time is in either exit state.
\end{lemma}

\begin{lemma}[Strict Guard-Checking Machine]
\lean{Turing.MachineFormat.exists_canonGuardFrontStrict}
\label{Turing.MachineFormat.exists_canonGuardFrontStrict}
\leanok
\uses{Turing.Cfg.mapState, Turing.Cfg.ofWords, Turing.FinTM, Turing.FinTM.bufferTape, Turing.MachineFormat, Turing.MachineFormat.canonGuardBlock, Turing.MachineFormat.canonHeader, Turing.MachineFormat.canonReadyWords, Turing.MachineFormat.canonSeam_avoids, Turing.MachineFormat.exists_canonGuardReady, Turing.MachineFormat.exists_canonReadySeam, Turing.MultiTapeTM.runFrom, Turing.codeBitsNat, Turing.pairDecode, Turing.seamCompTM, Turing.seamCompTM_run_ofCfg}
\difficulty{5}
For every machine format $F$, there exist a finite-state Turing machine $C$ over the
Boolean alphabet, with at least four work tapes, an entry state, an injective map from
$\{0,1\}$ to its states, and a constant $a\in\mathbb N$, such that the following
holds. Whenever bit strings $x,bits,rest$ satisfy that the aligned pair parser applied
to $x$ returns $(bits,rest)$, there is a time $t\le a(|x|+1)$ at which, starting on
input $x$ in the entry state with all work tapes blank, the machine is in the exit state
indexed by whether $B_F(\operatorname{val}(bits)+1)\le |rest|$, where
$B_F=\max(\operatorname{guard}(F),(4k+5)3^{k+1})$ and $\operatorname{val}(bits)$ is
the least-significant-bit-first binary value of $bits$. At that time its output is
empty, its input-head position is $2|bits|+3$ if the inequality holds and $|x|+1$
otherwise, and every work tape except tape $1$ contains its prescribed ready word with
its head at cell zero; moreover, no earlier time $u<t$ is in either exit state.
\end{lemma}

\begin{lemma}[Simulation with blank auxiliary tapes]
\lean{Turing.MachineFormat.canonBlankLift}
\label{Turing.MachineFormat.canonBlankLift}
\leanok
\uses{Turing.Cfg.ofWords, Turing.FinTM.bufferTape, Turing.MachineFormat.canonEmbedWords, Turing.MultiTapeTM, Turing.MultiTapeTM.runFrom, Turing.embedEmitCfg, Turing.embedEmitTM, Turing.embedEmitTM_runFrom}
\difficulty{3}
Let $m,k\in\mathbb{N}$, let $S$ be a type, let $M$ be a $m$-work-tape Turing machine
over the Boolean alphabet with states in $S$, and let $bank$ be an injection from the
$m$ work-tape indices into the $k$ work-tape indices. For every state $q\in S$, Boolean
input word $x$, and time $t\in\mathbb{N}$, running the machine obtained by embedding $M$
into the $k$-tape machine, from the configuration with state $q$ and all tapes blank,
gives the same configuration as embedding the configuration reached by running $M$ for
$t$ steps from that all-blank configuration; the additional tapes are blank, their heads
are at zero, and the output is unchanged.
\end{lemma}

\begin{lemma}[A machine for checking the canonical front]
\lean{Turing.MachineFormat.exists_canonCheckedFront}
\label{Turing.MachineFormat.exists_canonCheckedFront}
\leanok
\uses{Turing.Cfg.mapState, Turing.Cfg.ofWords, Turing.CodeFormat.decode, Turing.FinTM, Turing.FinTM.bufferTape, Turing.MachineFormat, Turing.MachineFormat.canonBlankLift, Turing.MachineFormat.canonEmbedWords, Turing.MachineFormat.canonGuardBlock, Turing.MachineFormat.canonHeaderOK, Turing.MachineFormat.canonHeaderOK_iff, Turing.MachineFormat.canonReadyWords, Turing.MachineFormat.canonSeam_avoids, Turing.MachineFormat.decode_of_canonGuard, Turing.MachineFormat.decode_of_canonHeader, Turing.MachineFormat.exists_canonGuardFrontStrict, Turing.MachineFormat.exists_canonHeaderBranch, Turing.MultiTapeTM.runFrom, Turing.codeBitsNat, Turing.embedEmitCfg_selected_pos, Turing.embedEmitCfg_selected_tape, Turing.embedEmitTM, Turing.embedEmitTM_frame, Turing.pairDecode, Turing.seamCompTM, Turing.seamCompTM_runFrom_left, Turing.seamCompTM_run_ofCfg}
\difficulty{7}
For every machine format $F$, there exist a finite Boolean-alphabet Turing machine $C$,
an entry state, an injective assignment of exit states to the three values in
$\{\mathrm{none},\mathrm{some\ false},\mathrm{some\ true}\}$, and a natural number $a$,
with at least four work tapes, such that the following holds for every Boolean input
$x$. At some time $t\le a(|x|+1)$, starting with all work tapes blank, the machine is in
the assigned exit state for a verdict, has produced no output, and has not previously
visited any exit state. If the verdict is not $\mathrm{some\ true}$, then decoding $x$
returns the fallback code. If the verdict is $\mathrm{some\ true}$, then $x$ parses as
$bits$ and $rest$, the binary representation of the value of $bits$ is exactly $bits$,
and $\max\{guard,(4k+5)3^{k+1}\}(v+1)\le |rest|$, where $k$ is the work-tape count of
$F$ and $v$ is the binary value of $bits$. At that time the input-head position is
$2|bits|+3$; every work tape other than tape $1$ is at position zero and contains the
canonical count word for $bits$ on tape $0$, the canonical offset word for $bits$ on
tape $2$, and the empty word on all remaining tapes.
\end{lemma}

\begin{lemma}[Dispatching an Exit State]
\lean{Turing.MachineFormat.canonExitDispatch}
\label{Turing.MachineFormat.canonExitDispatch}
\leanok
\uses{Turing.Cfg.mapState, Turing.Cfg.ofWords, Turing.MachineFormat.canonBoolPopTM, Turing.MultiTapeTM, Turing.MultiTapeTM.runFrom, Turing.phaseHostTM, Turing.phaseHostTM_run}
\difficulty{6}
Let $k>0$, let $S$ be a finite type with decidable equality, and let $M$ be a
$k$-work-tape Turing machine over Boolean symbols and states $S$. Choose a starting
state $s$, and an injective map $e$ from optional Boolean values to $S$. Then there are
a finite state type $R$ with decidable equality, a $k$-work-tape machine $N$ over
Boolean symbols and states $R$, an entry state in $R$, and an injective map $\ell$ from
bits to $R$ with the following property. For any Boolean input $x$, time $t$, and
optional bit $v$, suppose the run of $M$ from the configuration with state $s$ and empty
work tapes reaches state $e(v)$ at time $t$, and does not reach any state $e(q)$ at any
earlier time. Then the run of $N$ from its entry state with empty work tapes, on input
$x$, reaches at time $t+1$ the configuration obtained from that run of $M$ by replacing
its state with $\ell(1)$ when $v$ is the present bit $1$, and with $\ell(0)$ otherwise.
Moreover, before or at time $t+1$, the run of $N$ never reaches any state in the image
of $\ell$.
\end{lemma}

\begin{lemma}[Canonical parser front end]
\lean{Turing.MachineFormat.exists_canonParseFront}
\label{Turing.MachineFormat.exists_canonParseFront}
\leanok
\uses{Turing.Cfg.mapState, Turing.Cfg.ofWords, Turing.CodeFormat.decode, Turing.FinTM, Turing.FinTM.bufferTape, Turing.MachineFormat, Turing.MachineFormat.canonExitDispatch, Turing.MachineFormat.canonGuardBlock, Turing.MachineFormat.canonReadyWords, Turing.MachineFormat.exists_canonCheckedFront, Turing.MultiTapeTM.runFrom, Turing.codeBitsNat, Turing.pairDecode}
\difficulty{4}
For every machine format $F$, there exist a finite Turing machine $C$ over the Boolean
alphabet, an entry state $e$ of $C$, an injective map $q$ from Boolean verdicts to
states of $C$, and a natural number $a$, with at least four work tapes. For every
Boolean word $x$, the run from state $e$ on input $x$, with all work tapes blank, all
work-tape heads at cell zero, and empty output, reaches state $q(v)$ after some time
$t\le a(|x|+1)$, for a verdict $v$; no state $q(u)$ is reached at any time $u<t$. The
output remains empty at time $t$. If $v$ is false, decoding $x$ in format $F$ gives its
fallback code. If $v$ is true, there are bit strings $b,r$ such that pair decoding of
$x$ yields $(b,r)$, $b$ is the canonical binary representation of the natural number it
encodes, and $r$ has length at least $\max\{F.\mathrm{guard},(4C.k+5)3^{C.k+1}\}$ times
one more than that number. At time $t$ the input head is at position $2|b|+3$, and on
every work tape other than tape $1$, the tape contents store the canonical ready word
determined by $b$ and the tape head is at cell zero.
\end{lemma}

\begin{definition}[Canonical machine-code grammar]
\lean{Turing.MachineFormat.canonGrammar}
\label{Turing.MachineFormat.canonGrammar}
\leanok
\uses{Turing.codeSkipFin, Turing.codeSkipTableK}
Given natural numbers $k,n$ and a bit string $r$, the canonical grammar reader first
consumes a unary index smaller than $n+1$, then consumes a table for a machine with
$n+1$ states and $k$ work tapes. It succeeds exactly when every bit left over is $1$,
returning that remaining suffix; otherwise it fails.
\end{definition}

\begin{lemma}[Agreement of the scan and canonical grammar checks]
\lean{Turing.MachineFormat.scan_eq_canonGrammar}
\label{Turing.MachineFormat.scan_eq_canonGrammar}
\leanok
\uses{Turing.CodeFormat.scan, Turing.MachineFormat, Turing.MachineFormat.canonGrammar, Turing.MachineFormat.canonGuardBlock, Turing.codeBitsNat, Turing.codeSkipFin, Turing.pairDecode}
\difficulty{4}
Let $F$ be a machine format, and let $\alpha$, $b$, and $r$ be finite bit strings.
Suppose the aligned pair parser applied to $\alpha$ returns $(b,r)$, the string $b$ is
the canonical least-significant-bit-first binary representation of its value $n$, and
$\max\{g,(4k+5)3^{k+1}\}(n+1)\le |r|$, where $g$ is the format's guard and $k$ is its
number of work tapes. Then scanning $\alpha$ returns the same result as the following
check on $r$: consume a unary index smaller than $n+1$, skip the transition table for
$n+1$ states and $k$ work tapes, and accept exactly when every bit in the remaining
suffix is $1$, returning that suffix.
\end{lemma}

\begin{lemma}[Canonical Grammar Verdict]
\lean{Turing.MachineFormat.canonGrammar_verdict}
\label{Turing.MachineFormat.canonGrammar_verdict}
\leanok
\uses{Turing.CodeFormat.decode, Turing.CodeFormat.decode_of_parseFull_none, Turing.CodeFormat.scan, Turing.CodeFormat.scan_full, Turing.MachineFormat, Turing.MachineFormat.canonGrammar, Turing.MachineFormat.canonGuardBlock, Turing.MachineFormat.scan_eq_canonGrammar, Turing.codeBitsNat, Turing.pairDecode}
\difficulty{3}
Let $F$ be a machine format, and let $\alpha$, $bits$, and $rest$ be Boolean words.
Suppose parsing $\alpha$ as a pair gives $(bits,rest)$, the word $bits$ is the canonical
little-endian binary representation of its value $n$, and the remaining word has length
at least $B(n+1)$, where $B$ is the maximum of the format's guard constant and
$(4k+5)3^{k+1}$, with $k$ the format's work-tape count. If the grammar check on $n$ and
$rest$ fails, decoding $\alpha$ gives the format's fallback code. If the grammar check
succeeds, scanning $\alpha$ returns some suffix.
\end{lemma}

\begin{definition}[Canonical trailing-zero strip]
\lean{Turing.MachineFormat.canonStrip}
\label{Turing.MachineFormat.canonStrip}
\leanok
For a bit string, the canonical strip removes all trailing $1$ bits, leaving the prefix
ending at its last $0$, or the empty string if there is no $0$.
\end{definition}

\begin{lemma}[Removing a trailing run of ones]
\lean{Turing.MachineFormat.canonStrip_append}
\label{Turing.MachineFormat.canonStrip_append}
\leanok
\uses{Turing.MachineFormat.canonStrip}
\difficulty{3}
For Boolean words $u$ and $\mathrm{tail}$, if every bit of $\mathrm{tail}$ is $1$, then
removing all trailing $1$s from the word formed by concatenating $u$, the bit $0$, and
$\mathrm{tail}$ leaves exactly $u$ followed by $0$.
\end{lemma}

\begin{lemma}[Final Bit of a Flattened List]
\lean{Turing.MachineFormat.canonFlatMap_last}
\label{Turing.MachineFormat.canonFlatMap_last}
\leanok
\difficulty{4}
Let $L$ be a nonempty list, and let $f$ assign a Boolean word to each element of $L$. If
every word $f(a)$ for $a\in L$ ends in the bit $0$, then the concatenation of those
words in list order also ends in $0$.
\end{lemma}

\begin{lemma}[Final delimiter of an action encoding]
\lean{Turing.MachineFormat.actionBitsK_last}
\label{Turing.MachineFormat.actionBitsK_last}
\leanok
\uses{Turing.Action, Turing.actionBitsK, Turing.optBoolBits, Turing.optOptBoolBits, Turing.optStateBits, Turing.signBits, Turing.unaryFin}
\difficulty{4}
Let $k,n\in\mathbb{N}$, let $F$ be a machine format, and let $a$ be an action for a
$k$-work-tape machine with binary tape alphabet and state space $\{0,\ldots,n\}$. The
bit string encoding $a$ by the specified serialization of its input move, work-tape
instructions, output, and optional successor state ends in the bit $0$.
\end{lemma}

\begin{lemma}[Final zero in the machine serialization]
\lean{Turing.MachineFormat.serializeK_last}
\label{Turing.MachineFormat.serializeK_last}
\leanok
\uses{Turing.MachineFormat.actionBitsK_last, Turing.MachineFormat.canonFlatMap_last, Turing.MultiTapeTM, Turing.StepBlock.length_workReadsK, Turing.actionBitsK, Turing.pairEncode, Turing.serializeK, Turing.tableBitsK, Turing.unaryFin, Turing.workReadsK}
\difficulty{5}
For natural numbers $k$ and $n$, let $M$ be a deterministic Turing machine with $k$ work
tapes, binary tape alphabet, and state set $	ext{Fin}(n+1)$. Its serialization can be
written as a finite bit string followed by a final $0$.
\end{lemma}

\begin{lemma}[Accepted scans leave only ones]
\lean{Turing.MachineFormat.canonScan_all}
\label{Turing.MachineFormat.canonScan_all}
\leanok
\uses{Turing.CodeFormat.scan, Turing.MachineFormat}
\difficulty{5}
Let $F$ be a code format and let $\alpha$ and $tail$ be Boolean words. If scanning
$\alpha$ in format $F$ succeeds with returned suffix $tail$, then every bit of $tail$ is
$1$.
\end{lemma}

\begin{lemma}[Canonical prefix from a successful scan]
\lean{Turing.MachineFormat.canonStrip_of_scan}
\label{Turing.MachineFormat.canonStrip_of_scan}
\leanok
\uses{Turing.CodeFormat.decode, Turing.CodeFormat.parseFull_sound, Turing.CodeFormat.scan, Turing.CodeFormat.scan_full, Turing.MachineFormat, Turing.MachineFormat.canonScan_all, Turing.MachineFormat.canonStrip, Turing.MachineFormat.canonStrip_append, Turing.MachineFormat.serializeK_last}
\difficulty{4}
Let $F$ be a machine format, and let $\alpha$ and $\mathrm{tail}$ be Boolean words. If
scanning $\alpha$ succeeds and returns $\mathrm{tail}$, then removing all trailing $1$s
from $\alpha$ gives the serialization of the code obtained by decoding $\alpha$.
\end{lemma}

\begin{lemma}[Decomposition after removing trailing ones]
\lean{Turing.MachineFormat.canonStrip_split}
\label{Turing.MachineFormat.canonStrip_split}
\leanok
\uses{Turing.MachineFormat.canonStrip}
\difficulty{4}
For every Boolean word $w$, there is a natural number $r$ such that $w$ is the
concatenation of the word obtained by removing all trailing $1$s from $w$ and a string
of $r$ copies of $1$. The resulting prefix is either empty or ends in $0$.
\end{lemma}

\begin{lemma}[Boundary properties of trailing-one removal]
\lean{Turing.MachineFormat.canonStrip_boundary}
\label{Turing.MachineFormat.canonStrip_boundary}
\leanok
\uses{Turing.MachineFormat.canonStrip, Turing.MachineFormat.canonStrip_split}
\difficulty{4}
For every Boolean word $w$, let $v$ be the word obtained by removing all trailing $1$s
from $w$. Then $|v|\le |w|$ and the prefix of $w$ of length $|v|$ is $v$. Every bit of
$w$ at an index $r<|w|$ with $|v|\le r$ is $1$, and whenever $r<|w|$ and $|v|=r+1$, the
bit of $w$ at index $r$ is $0$.
\end{lemma}

\begin{definition}[Canonical trimming machine]
\lean{Turing.MachineFormat.canonTrimTM}
\label{Turing.MachineFormat.canonTrimTM}
\leanok
\uses{Turing.MultiTapeTM, Turing.ScanPhase, Turing.tapeAction}
For a natural number $k$ and a work-tape index $i<k$, define a deterministic machine
with Boolean tape symbols and states recording a forward scan, a backward scan with a
Boolean flag, or a completed scan with a Boolean flag. It starts in the forward state
with flag false; while scanning forward it moves right on a nonblank symbol and moves
left into the backward phase on a blank, and while scanning backward it moves left on a
nonblank symbol, setting the flag to the conjunction of its current value and the
scanned symbol, and moves right into the completed phase on a blank. In the backward
phase, when the flag and scanned symbol are both true, it also prescribes writing a
blank; in the completed phase it stays in place and remains completed.
\end{definition}

\begin{lemma}[Erasing the Final Tape Cell]
\lean{Turing.MachineFormat.canonEraseLast}
\label{Turing.MachineFormat.canonEraseLast}
\leanok
\uses{Turing.FinTM.bufferTape, Turing.FinTM.bufferTape_append}
\difficulty{3}
For every Boolean word $w$ and bit $b$, take the tape that stores the word formed by
appending $b$ to $w$, and erase its cell at index $|w|$. The resulting tape content is
exactly the tape that stores $w$.
\end{lemma}

\begin{lemma}[Run of the tape trimming machine]
\lean{Turing.MachineFormat.canonTrim_run}
\label{Turing.MachineFormat.canonTrim_run}
\leanok
\uses{Turing.Cfg, Turing.Cfg.ofWords, Turing.Cfg.workTapeSymbols, Turing.FinTM.bufferTape, Turing.FinTM.bufferTape_nat, Turing.MachineFormat.canonEraseLast, Turing.MachineFormat.canonStrip, Turing.MachineFormat.canonStrip_boundary, Turing.MachineFormat.canonTrimTM, Turing.MultiTapeTM.runFrom, Turing.MultiTapeTM.step, Turing.ScanPhase, Turing.catalogTrace, Turing.catalog_trace_run, Turing.moveInputPos_zero, Turing.step_of_tapeAction_move, Turing.step_of_tapeAction_write}
\difficulty{7}
Let $k$ be a natural number, let $i$ be one of $k$ work tapes, let $x$ be a Boolean
input word, and let $\mathrm{words}$ assign a Boolean word to each work tape. Starting
with the machine in its forward-pass state on input $x$ and with these words on the
tapes, after exactly $2|\mathrm{words}(i)|+2$ steps it is in its completed state, with
tape $i$ replaced by the word obtained by removing its trailing $1$s and every other
tape unchanged. It does not enter that completed state at any earlier time.
\end{lemma}

\begin{lemma}[Polynomial-time canonical stripping]
\lean{Turing.MachineFormat.exists_canonStripTM}
\label{Turing.MachineFormat.exists_canonStripTM}
\leanok
\uses{Turing.Action.apply, Turing.Cfg, Turing.Cfg.mapState, Turing.Cfg.ofWords, Turing.FinTM, Turing.FinTM.ComputesFunInTime, Turing.FinTM.ComputesInTime, Turing.FinTM.ComputesInTime.mono, Turing.FinTM.bufferTape, Turing.FinTM.bufferTape_nil, Turing.FinTM.computesFunInTime_id, Turing.FinTM.computesInTime_iff, Turing.FinTM.exists_emitCallTM, Turing.MachineFormat.canonStrip, Turing.MachineFormat.canonStrip_boundary, Turing.MachineFormat.canonTrimTM, Turing.MachineFormat.canonTrim_run, Turing.MultiTapeTM, Turing.MultiTapeTM.runFrom, Turing.MultiTapeTM.runFrom_succ_eq_step, Turing.MultiTapeTM.step, Turing.ScanPhase, Turing.SweepPhase, Turing.initCfg_ofWords, Turing.inputCopyTM, Turing.inputCopyTM_run_ofCfg, Turing.seamCompTM, Turing.seamCompTM_firstReturn_ofCfg, Turing.seamCompTM_run_ofCfg, Turing.seamReleaseTM, Turing.seamReleaseTM_runFrom_pos, Turing.stateWord}
\difficulty{7}
There are a finite Turing machine $M$ over the Boolean alphabet and a natural number $a$
such that, on every Boolean word $\alpha$, $M$ halts after $a(|\alpha|+1)^a$ steps with
output equal to the word obtained by deleting all trailing ones from $\alpha$.
\end{lemma}

\begin{lemma}[A machine for the grammar-checking stage]
\lean{Turing.MachineFormat.exists_canonGrammarStage}
\label{Turing.MachineFormat.exists_canonGrammarStage}
\leanok
\uses{Turing.Cfg, Turing.FinTM.bufferTape, Turing.MachineFormat, Turing.MachineFormat.canonGrammar, Turing.MachineFormat.canonGuardBlock, Turing.MultiTapeTM, Turing.MultiTapeTM.runFrom, Turing.codeBitsNat, Turing.pairDecode}
For every machine format $F$, there exist a natural number $m$, a finite state set $S$
with decidable equality, a deterministic Turing machine $G$ with $m$ work tapes over the
Boolean alphabet, an entry state, an injective assignment of exit states to Boolean
values, a natural number $a$, and a designated work-tape index $header$. For every input
$x$ that parses as $bits$ and $rest$, with $bits$ the canonical binary representation of
its value and $|rest|\ge \max\{guard,(4k+5)3^{k+1}\}(v+1)$, where $k$ is the work-tape
count of $F$ and $v$ is the value of $bits$, and for every configuration $d$ on input
$x$ whose state is the entry state, output is empty, input-head position is $2|bits|+3$,
designated tape contains $bits$ with its head at zero, and all other tapes are blank
with heads at zero, there is a time $t\le a(|x|+1)^a$ such that $G$ has empty output and
is in the exit state assigned to whether the grammar check succeeds. No exit state is
visited before time $t$.
\end{lemma}

\begin{lemma}[A bounded-time decider for decoding]
\lean{Turing.MachineFormat.exists_canonDecider}
\label{Turing.MachineFormat.exists_canonDecider}
\leanok
\uses{Turing.CodeFormat.decode, Turing.CodeFormat.scan, Turing.FinTM, Turing.FinTM.ComputesInTime, Turing.MachineFormat, Turing.MachineFormat.canonGrammar_verdict, Turing.MachineFormat.exists_canonGrammarStage, Turing.MachineFormat.exists_canonParseFront}
\difficulty{4}
For every machine format $F$, there exist a finite Turing machine $D$ over the Boolean
alphabet and a natural number $a$ such that for each Boolean string $\alpha$, some
Boolean value $b$ is output as the single-symbol word $[b]$ within $a(|\alpha|+1)^a$
steps. If $b$ is false, decoding $\alpha$ returns the fallback code; if $b$ is true, the
scan of $\alpha$ succeeds and returns a suffix.
\end{lemma}

\begin{theorem}[Polynomial-time canonical serialization]
\lean{Turing.MachineFormat.exists_polyCanonizer}
\label{Turing.MachineFormat.exists_polyCanonizer}
\leanok
\uses{Turing.CodeFormat.decode, Turing.FinTM, Turing.FinTM.ComputesFunInTime, Turing.FinTM.ComputesInTime.mono, Turing.FinTM.computesFunInTime_cond, Turing.FinTM.computesFunInTime_const, Turing.MachineFormat.canonStrip, Turing.MachineFormat.canonStrip_of_scan, Turing.MachineFormat.exists_canonDecider, Turing.MachineFormat.exists_canonStripTM}
\difficulty{6}
For every machine format $F$, there are a finite Turing machine $C$ over the Boolean
alphabet and a natural number $d$ such that, on every Boolean input $\alpha$, $C$ halts
within $d(|\alpha|+1)^d$ steps with output equal to the serialization, in format $F$, of
the code decoded from $\alpha$.
\end{theorem}

\begin{definition}[Summary clock]
\lean{Turing.MachineFormat.summaryClock}
\label{Turing.MachineFormat.summaryClock}
\leanok
\uses{Turing.Cfg, Turing.MultiTapeTM, Turing.MultiTapeTM.outputSymbol, Turing.MultiTapeTM.step, Turing.StepBlock.summaryEmit}
Let $M$ be a deterministic machine with $k$ work tapes, Boolean symbols, and state set
$Q$, and let $x$ be a finite Boolean input. Define a Boolean-valued summary from a
configuration $c$, an optional Boolean summary $s$, and a natural time bound. At time
zero its value is true exactly when $c$ is halted and $s$ is the Boolean value true; at
one more than time $t$, its value is the summary at time $t$ of the successor
configuration after one step of $M$, using as the updated summary the result of
combining $s$ with the optional symbol emitted by that step (retaining $s$ if no symbol
is emitted, setting the result to true when the first emitted symbol is true, and
setting it to false otherwise).
\end{definition}

\begin{lemma}[Summary clock correctness]
\lean{Turing.MachineFormat.summaryClock_correct}
\label{Turing.MachineFormat.summaryClock_correct}
\leanok
\uses{Turing.Cfg, Turing.MachineFormat.summaryClock, Turing.MultiTapeTM, Turing.MultiTapeTM.runFrom, Turing.MultiTapeTM.runFrom_succ_eq_step, Turing.MultiTapeTM.runFrom_zero, Turing.MultiTapeTM.step_output, Turing.StepBlock.outputSummary, Turing.StepBlock.outputSummary_append}
\difficulty{4}
Let $M$ be a multi-tape Turing machine over bits with state set $Q$, let $x$ be a finite
bit string, let $c$ be a configuration of $M$ on input $x$, and let $t$ be a natural
number. Starting the summary clock at $c$ with the summary of its current output gives
true exactly when the configuration reached after $t$ steps is halted and the summary of
its output is true.
\end{lemma}

\begin{lemma}[Summary clock recognizes acceptance]
\lean{Turing.MachineFormat.summaryClock_accepts}
\label{Turing.MachineFormat.summaryClock_accepts}
\leanok
\uses{Turing.FinTM, Turing.FinTM.ComputesInTime, Turing.FinTM.computesInTime_iff, Turing.MachineFormat.summaryClock, Turing.MachineFormat.summaryClock_correct, Turing.MultiTapeTM.initCfg, Turing.StepBlock.outputSummary, Turing.StepBlock.outputSummary_true}
\difficulty{3}
Let $M$ be a finite-state multi-tape Turing machine over the binary alphabet, let $x$ be
a finite bit string, and let $t$ be a natural number. The summary clock started on $x$
with no output summary equals true exactly when $M$ halts after $t$ steps with output
the one-bit word $1$ and uses some finite amount of work-tape space.
\end{lemma}

\begin{lemma}[Summary clock at time zero]
\lean{Turing.MachineFormat.summaryClock_zero}
\label{Turing.MachineFormat.summaryClock_zero}
\leanok
\uses{Turing.FinTM, Turing.MachineFormat.summaryClock, Turing.MultiTapeTM.initCfg}
\difficulty{1}
For every finite-state multi-tape Turing machine $M$ over bits and every finite bit
string $x$, the summary clock at time zero, started in its initial configuration with no
output summary, equals false.
\end{lemma}

\begin{definition}[Hosted verdict machine]
\lean{Turing.MachineFormat.h1HostedVerdictTM}
\label{Turing.MachineFormat.h1HostedVerdictTM}
\leanok
\uses{Turing.CounterLoopState, Turing.MultiTapeTM, Turing.StepBlock.Ph, Turing.StepBlock.hostedLoopTM, Turing.StepBlock.summaryTape, Turing.vhostBank}
For each natural number $k$, define a Boolean-symbol machine with $1+(k+2+1)$ work tapes
and states consisting of a counter-loop control state, a phase indexed by a key, and a
Boolean flag. Its initial state and transition behavior are those of the hosted loop
machine, except that in the designated completed or escape control states it stays in
place, halts, and emits a Boolean verdict: in the escape state this is true exactly when
the scanned cell on the designated summary tape contains true, and in the completed
state it is false.
\end{definition}

\begin{theorem}[Hosted summary effect]
\lean{Turing.MachineFormat.h1HostedVerdict_effect}
\label{Turing.MachineFormat.h1HostedVerdict_effect}
\leanok
\uses{Turing.Action.apply, Turing.Cfg, Turing.Cfg.init, Turing.FinTM.VirtualTag, Turing.FinTM.bufferTape, Turing.MachineFormat.h1HostedVerdictTM, Turing.MachineFormat.summaryClock, Turing.MachineFormat.summaryClock_correct, Turing.MultiTapeTM, Turing.MultiTapeTM.AgreeOn, Turing.MultiTapeTM.initCfg, Turing.MultiTapeTM.runFrom, Turing.MultiTapeTM.runFrom_eq_of_agreeOn, Turing.MultiTapeTM.runFrom_succ_eq_step', Turing.MultiTapeTM.step, Turing.StepBlock.Ph, Turing.StepBlock.Rep, Turing.StepBlock.exit, Turing.StepBlock.hostedLoopTM, Turing.StepBlock.hostedLoopTM_run, Turing.StepBlock.hostedLoop_read, Turing.StepBlock.outputSummary, Turing.counterLoopCfg, Turing.counterValue_bits, Turing.counter_cfg_eq, Turing.moveInputPos_zero, Turing.tableBitsK, Turing.vhostCfg}
\difficulty{6}
Let $M$ be a deterministic machine with $k$ work tapes and state set $\{0,\ldots,n\}$,
let $x,z$ be finite bit strings, and let $d$ be a configuration of the step-block
machine for $M$ on input $x$. Suppose $d$ is at the live entry state and correctly
represents $M$'s initial configuration together with its transition table and summary
tapes. For every $t\in\mathbb N$ and every input-head position $p$ among the $|z|+2$
positions, the hosted verdict machine started from the specified counter initialization
of $d$, with input-head position $p$, halts after some $T\le
t\bigl(50(|L|+1)+4\bigr)+2\ell(t)+3$ steps, where $L$ is the serialized transition table
and $\ell(t)$ is the length of the binary expansion of $t$; its final input-head
position is $p$ and its output is $d$'s output followed by $s_t$. Here $c_0$ is the
initial configuration of $M$ on $x$, $s_0$ is true exactly when $c_0$ is already halted,
and $s_{j+1}$ is the summary obtained by combining $s_j$ with the output symbol emitted
by the transition from $c_j$, while $c_{j+1}$ is the next configuration after that
transition.
\end{theorem}

\begin{theorem}[Hosted machine verdict]
\lean{Turing.MachineFormat.h1HostedVerdict_accepts}
\label{Turing.MachineFormat.h1HostedVerdict_accepts}
\leanok
\uses{Turing.Cfg, Turing.CodeFormat.decode, Turing.FinTM.ComputesInTime, Turing.FinTM.bufferTape, Turing.MachineFormat, Turing.MachineFormat.h1HostedVerdictTM, Turing.MachineFormat.h1HostedVerdict_effect, Turing.MachineFormat.summaryClock_accepts, Turing.MachineFormat.toFinTM, Turing.MultiTapeTM.initCfg, Turing.MultiTapeTM.runFrom, Turing.StepBlock.Ph, Turing.StepBlock.Rep, Turing.StepBlock.exit, Turing.counterLoopCfg, Turing.tableBitsK, Turing.vhostCfg}
\difficulty{1}
Let $F$ be a machine format, let $\alpha,x,z$ be finite bit strings, and let $d$ be a
configuration of the step-block machine for the machine denoted by decoding $\alpha$, on
input $x$. Suppose $d$ is at the live entry state and correctly represents that decoded
machine's initial configuration together with its transition table and summary tapes.
For every $t\in\mathbb N$ and every input-head position $p$ among the $|z|+2$ positions,
the hosted verdict machine, started with $d$, a counter encoding $t$, and input-head
position $p$, halts after some $T$ steps with $T\le
t\bigl(50(|L|+1)+4\bigr)+2\ell(t)+3$, where $L$ is the serialized transition table and
$\ell(t)$ is the length of the binary expansion of $t$. Its output is $d$'s output
followed by a bit $a$, and $a=1$ if and only if the decoded machine halts on $x$ within
$t$ steps with output exactly $[1]$.
\end{theorem}

\begin{theorem}[Uniform simulation of coded machines]
\lean{Turing.MachineFormat.exists_uniformSimulator}
\label{Turing.MachineFormat.exists_uniformSimulator}
\leanok
\uses{Turing.CodeFormat.decode, Turing.FinTM, Turing.FinTM.ComputesInTime, Turing.MachineFormat.toFinTM, Turing.pairEncode}
For every machine format $F$, there are a finite Turing machine $S$ over the Boolean
alphabet and a natural number $d$ such that, for every pair of Boolean words $\alpha,x$
and every $t\in\mathbb N$, on input encoding the pair consisting of the binary
representation of $t$ and $\alpha$, paired with $x$, the machine $S$ halts within
$d(|\alpha|+|x|+t+1)^d$ steps and outputs the single bit $1$ exactly when the machine
denoted by the code decoded from $\alpha$ halts on $x$ within $t$ steps with output
$[1]$; otherwise it outputs the single bit $0$.
\end{theorem}
